\subsection{\texorpdfstring{Properties of mappings \(\Theta_{a}\)}{Properties of mappings \textbackslash Theta\_\{a\}}}

Recall (VAR, R2, pp.~46 and~47, 11.1.3, d)) that if K is a compact subset of C, there exists a fundamental system of compact neighborhoods V of K which are pieces of C (that is, for every \(x \in V\) there exists a chart \((\mathrm{U}, \varphi, \mathbf{C})\) from C at x such that \(\varphi(\mathrm{U} \cap \mathrm{K})\) is an open subset of a closed half-space of C). We then denote by \(\partial V\) the boundary of the piece V equipped with the orientation induced by the orientation of C (VAR, R2, p.~47) and let dz denote the differential of the injection \(\partial V \to C\) .

PROPOSITION 1. --- Let a be an element of A, let U be an open neighborhood of Sp(a), and let \(f \in \mathcal{O}(\mathrm{U};\mathrm{A})\). Let V be a compact neighborhood of Sp(a) contained in U and such that V is a piece of C.

Then \(z \mapsto f(z)(z - a)^{-1}\) is continuous on \(\partial\mathrm{V}\); the differential form \(f(z)(z - a)^{-1}dz\) is integrable on \(\partial\mathrm{V}\), and one has

\[
\Theta_ {a} (f) = \frac {1}{2 i \pi} \int_ {\partial \mathrm{V}} f (z) (z - a) ^ {- 1} d z.
\]

Let \(h\) be a map of \(\mathbf{C}\) into \(\mathbf{C}\) infinitely differentiable, equal to 1 in a neighborhood of \(\mathrm{Sp}(a)\) and with compact support contained in the interior of \(V\) (lemma 1 of I, p.~52). Let \(u\) be a map of \(\mathbf{C}\) to A such that \((h, u)\) is adapted to \(a\) (cf. example 3 of I, p. 53). The associated differential form is \(\omega = du \wedge dz\). It follows that \(f\omega = fdu \wedge dz = d(fu dz)\) since \(f\) is holomorphic. Moreover, \(u(z) = (z - a)^{-1}\) on the boundary of \(V\) . Furthermore, the differential form \(fu dz\) is of class \(C^1\) on U. Therefore,

\[
2 i \pi \Theta_ {a} (f) = \int_ {\mathrm{V}} d (f u d z) = \int_ {\partial \mathrm{V}} f u d z = \int_ {\partial \mathrm{V}} f (z) (z - a) ^ {- 1} d z
\]

by formula (5) and Stokes' formula for the piece V (VAR, R2, p.~47, 11.2.3).

\[
\text { COROLLARY. } - \text {   Let   } a \in \mathrm{A}. \text {   One has   } \Theta_ {a} (1) = 1.
\]

Let \(\mathrm{R} > \varrho(a)\) be a real number. Let V be the closed disk with center 0 and radius R, so that \(\mathrm{Sp}(a) \subset \mathring{\mathrm{V}}\) . This is a piece of C whose boundary \(\partial\mathrm{V}\) is the circle with center 0 and radius R. For \(z \in \mathbf{C} - \mathring{\mathrm{V}}\) , we have the formula \((z - a)^{-1} = z^{-1}(1 - z^{-1}a)^{-1} = \sum_{j=1}^{+\infty} z^{-j} a^{j-1}\) . The series converges uniformly for \(z \in \partial V\) . We therefore have

\[
\Theta_ {a} (1) = \frac {1}{2 i \pi} \sum_ {j = 1} ^ {+ \infty} a ^ {j - 1} \int_ {\partial \mathrm{V}} z ^ {- j} d z = 1
\]

since

\[
\int_ {\partial \mathrm{V}} z ^ {j} d z = 0
\]

for every integer \(j \neq -1\) and

\[
\int_ {\partial \mathrm{V}} z ^ {- 1} d z = 2 i \pi
\]

(VAR, R2, p.~44, 10.4.5, and p.~47, 11.2.1, example).

Lemma 8. --- Let U be an open subset of \(C^{n}\) . Let \(\omega_{1}\) be a continuous differential form of degree n in U with compact support and values in A (resp. let \(\omega_{2}\) be a continuous differential form of degree 2 in C with compact support and values in C). Let us denote by \(\pi_{1}\) and \(\pi_{2}\) the canonical projections of U × C onto U and C. The differential form \(\pi_{1}^{*}\omega_{1} \wedge \pi_{2}^{*}\omega_{2}\) on U × C is continuous with compact support and takes values in A. We have

\[
\int_ {\mathrm{U} \times \mathbf {C}} \pi_ {1} ^ {*} \omega_ {1} \wedge \pi_ {2} ^ {*} \omega_ {2} = \left(\int_ {\mathbf {C}} \omega_ {2}\right) \left(\int_ {\mathrm{U}} \omega_ {1}\right).
\]

Let \(\mu_{n}\) be the Lebesgue measure on \(C^{n}\) and \(\mu_{1}\) the Lebesgue measure on C. There exists \(\psi_{1} \in \mathcal{K}(U;A)\) and \(\psi_{2} \in \mathcal{K}(C)\) such that the vector measure associated with \(\omega_{1}\) is equal to \(\psi_{1} \cdot \mu_{n}\) , and the vector measure associated with \(\omega_{2}\) is equal to \(\psi_{2} \cdot \mu_{1}\) . The vector measure associated with the differential form \(\pi_{1}^{*}\omega_{1} \wedge \pi_{2}^{*}\omega_{2}\) is \((\psi_{1} \otimes \psi_{2}) \cdot \mu_{n} \otimes \mu_{1}\) .

Let \(\ell\) be a continuous linear form on A. By INT, VI, §2, no. 2, def. 2 and the definition of the product measure (INT, III, §4, no. 1, def. 1), it follows

\[
\begin{array}{r l} & {\ell \Big (\int_ {\mathrm{U} \times \mathbf {C}} \pi_ {1} ^ {*} \omega_ {1} \wedge \pi_ {2} ^ {*} \omega_ {2} \Big) = \int_ {\mathrm{U} \times \mathbf {C}} \ell \circ (\psi_ {1} \otimes \psi_ {2}) \mu_ {n} \otimes \mu_ {1}} \\ & {\qquad = \int_ {\mathrm{U} \times \mathbf {C}} \psi_ {2} (z) \ell (\psi_ {1} (x)) d \mu_ {n} (x) d \mu_ {1} (z)} \\ & {\qquad = \Big (\int_ {\mathbf {C}} \psi_ {2} (z) d \mu_ {1} (z) \Big) \Big (\int_ {\mathrm{U}} \ell (\psi_ {1} (x)) d \mu_ {n} (x) \Big)} \\ & {\qquad = \Big (\int_ {\mathbf {C}} \psi_ {2} \mu_ {1} \Big) \ell \Big (\int_ {\mathrm{U}} \psi_ {1} \mu_ {n} \Big),} \end{array}
\]

whence the result (INT VI, loc. cit.).

Lemma 9. --- Let \(\boldsymbol{a} = (a_1, \ldots, a_n) \in \mathbb{A}^n\) . Let \(p \in \mathbf{N}\) and \(\boldsymbol{a}' = (a_{n+1}, \ldots, a_{n+p}) \in \mathbb{A}^p\) . Then we have \(\Theta_{(\boldsymbol{a}, \boldsymbol{a}')} \circ \pi_{n,n+p}^* = \Theta_{\boldsymbol{a}}\) . In particular, we have \(\Theta_{\boldsymbol{a}}(1) = 1\) .

As we have \(\pi_{n,n+p} = \pi_{n,n+1} \circ \cdots \circ \pi_{n+p-1,n+p}\) , it suffices to prove the first assertion when \(p = 1\) , which we now assume. We simply denote \(\pi = \pi_{n,n+1}\) . It then suffices to show that, for every open neighborhood U of \(\mathrm{Sp}^n(\boldsymbol{a})\) , and every function \(f \in \mathcal{O}(\mathrm{U};\mathrm{A})\) , one has \(\Theta_{(\boldsymbol{a},a_{n+1})}(f \circ \pi) = \Theta_{\boldsymbol{a}}(f)\) . Denote by \(g = f \circ \pi\) . Let \(h\) (resp. \(h'\)) be a mapping from \(\mathbf{C}^n\) into \(\mathbf{C}\) (resp. from \(\mathbf{C}\) into \(\mathbf{C}\)), infinitely differentiable, equal to 1 in a neighborhood of \(\mathrm{Sp}^n(\boldsymbol{a})\) (resp. of \(\mathrm{Sp}(\boldsymbol{a}')\)), with compact support contained in U (resp. in \(\mathbf{C}\) ). There exist mappings ( \(u_1, \ldots, u_n\) ) from \(\mathbf{C}^n\) into A, infinitely differentiable, such that the sequence ( \(h, u_1, \ldots, u_n\) ) is adapted to \(\boldsymbol{a}\) (lemma 3 of I, p.~54), and an infinitely differentiable mapping \(u_{n+1}\) of \(\mathbf{C}\) into A such that the pair ( \(h', u_{n+1} \) ) is adapted to \(a_{n+1}\) (loc. cit.)

For \(z \in C^{n}\) and \(z_{n+1} \in C\) , let \(h''(z, z_{n+1}) = h(z)h'(z_{n+1})\) and \(u_{n+1}''(z, z_{n+1}) = h(z)u_{n+1}(z_{n+1})\) . The functions \(h''\) and \(u_{n+1}''\) are infinitely differentiable in \(C^{n+1}\) . The function \(h''\) is equal to 1 in a neighborhood of \(\mathrm{Sp}^{n+1}(a, a_{n+1})\) , and has compact support contained in \(U \times C\) . For every \(\boldsymbol{w} = (z, z_{n+1}) \in \mathbf{C}^{n+1}\) , one has

\[
\begin{array}{c} (z _ {1} - a _ {1}) (u _ {1} \circ \pi) (\boldsymbol {w}) + \dots + (z _ {n} - a _ {n}) (u _ {n} \circ \pi) (\boldsymbol {w}) \\ + (z _ {n + 1} - a _ {n + 1}) u _ {n + 1} ^ {\prime \prime} (\boldsymbol {w}) = 1 - h (\boldsymbol {z}) + h (\boldsymbol {z}) (1 - h ^ {\prime} (z _ {n + 1})) \\ = 1 - h ^ {\prime \prime} (\boldsymbol {w}), \end{array}
\]

which shows that the sequence \((h''_{1}, u_{1} \circ \pi, \ldots, u_{n} \circ \pi, u_{n+1}')\) is adapted to \((\boldsymbol{a}, a_{n+1})\) . Let \(\omega\) be the associated differential form.

The differential form \(du_{1} \wedge dz_{1} \wedge \cdots \wedge du_{n} \wedge dz_{n} \wedge dh\) on \(C^{n}\) has degree \(2n + 1\) , hence is zero. Consequently,

\[
\begin{array}{l} \omega = d (u _ {1} \circ \pi) \wedge d z _ {1} \wedge \dots \wedge d (u _ {n} \circ \pi) \wedge d z _ {n} \wedge d u _ {n + 1} ^ {\prime \prime} \wedge d z _ {n + 1} \\ = (h \circ \pi) d (u _ {1} \circ \pi) \wedge d z _ {1} \wedge \dots \wedge d (u _ {n} \circ \pi) \wedge d z _ {n} \wedge d u _ {n + 1} \wedge d z _ {n + 1}. \end{array}
\]

Since \(g = f \circ \pi\) , formula (5) and lemma 8 imply

\[
\begin{array}{c} \Theta_ {\boldsymbol {a}, \boldsymbol {a} ^ {\prime}} (g) = \frac {(n + 1) !}{(2 i \pi) ^ {n + 1}} \int_ {\mathrm{U} \times \mathbf {C}} g \omega = \frac {(n + 1) !}{(2 i \pi) ^ {n + 1}} \Big (\int_ {\mathbf {C}} d u _ {n + 1} \wedge d z _ {n + 1} \Big) \\ \qquad \qquad \times \Big (\int_ {\mathrm{U}} f h   d u _ {1} \wedge d z _ {1} \wedge \dots \wedge d u _ {n} \wedge d z _ {n} \Big). \end{array}
\]

On the one hand, we have

\[
\int_ {\mathbf {C}} d u _ {n + 1} \wedge d z _ {n + 1} = 2 i \pi \Theta_ {a _ {n + 1}} ^ {\mathbf {C}} (1) = 2 i \pi \cdot 1
\]

by corollary 5. On the other hand, part c) of lemma 4 of I, p.~54 and the fact that the integral of a closed form is zero (VAR, R2, p.~48, 11.2.4) imply

\[
\begin{array}{c} (n + 1) \int_ {\mathrm{U}} f h d u _ {1} \wedge d z _ {1} \wedge \dots \wedge d u _ {n} \wedge d z _ {n} = \\ \int_ {\mathrm{U}} f   d u _ {1} \wedge d z _ {1} \wedge \dots \wedge d u _ {n} \wedge d z _ {n} = \frac {(2 i \pi) ^ {n}}{n !} \Theta_ {\boldsymbol {a}} (f). \end{array}
\]

Thus one obtains

\[
\Theta_ {\boldsymbol {a}, \boldsymbol {a} ^ {\prime}} (g) = \frac {(n + 1) !}{(2 i \pi) ^ {n + 1}} \times \frac {(2 i \pi) ^ {n}}{(n + 1) !} \Theta_ {\boldsymbol {a}} (f) \times 2 i \pi = \Theta_ {\boldsymbol {a}} (f).
\]

Finally, the formula \(\Theta_{\boldsymbol{a}}(1)=1\) follows from the foregoing and from the corollary of Prop. 1.

Lemma 10. --- Let \(\boldsymbol{a} \in \mathrm{A}^n\) . Let \(g\) be a polynomial function on \(\mathbf{C}^n\) with coefficients in A and \(f \in \mathcal{O}(\mathrm{Sp}^n(\boldsymbol{a});\mathrm{A})\) . We have \(\Theta_{\boldsymbol{a}}(gf) = g(\boldsymbol{a})\Theta_{\boldsymbol{a}}(f)\) . In particular, we have \(\Theta_{\boldsymbol{a}}(g) = g(\boldsymbol{a})\) .

By Lemma 9, it suffices to prove the first assertion.

We denote by \(z_{1},\ldots,z_{n}\) the coordinate functions on \(C^{n}\). Since the map \(\Theta_{a}\) is A-linear, it suffices to prove the assertion of the lemma when \(g=z_{1}^{e_{1}}\cdots z_{n}^{e_{n}}\) , where \((e_{1},\ldots,e_{n})\in\mathbf{N}^{n}\). Proceeding by induction on \(e_{1}+\cdots+e_{n}\) , we reduce to the case where there exists an integer i such that \(1\leqslant i\leqslant n\) and \(g=z_{i}\) .

Let U be an open neighborhood of \(\mathrm{Sp}^{n}(\boldsymbol{a})\) . Let \((h,u_{1},\ldots,u_{n})\) be a sequence adapted to \(\boldsymbol{a}\) such that the support of h is contained in U (Lemma 1 of I, p.~52 and Lemma 3 of I, p.~54), and let \(\omega\) be the associated differential form. By lemma 4, a) of I, p.~54, there exists a differential form \(\beta\) such that

\[
(z _ {i} - a _ {i}) \omega = d (h \beta \wedge d z _ {1} \wedge \dots \wedge d z _ {n}).
\]

Consequently, for every function \(f \in \mathcal{O}(\mathrm{U};\mathrm{A})\) , one has

\[
\left(z _ {i} - a _ {i}\right) f \omega = f d \left(h \beta \wedge d z _ {1} \wedge \dots \wedge d z _ {n}\right) = d \left(f h \beta \wedge d z _ {1} \wedge \dots \wedge d z _ {n}\right)
\]

since f is holomorphic, so that \(df \wedge dz_{1} \wedge \cdots \wedge dz_{n} = 0\). Applying Stokes' formula (VAR, R2, p.~48, 11.2.4), we obtain \(\int_{\mathrm{U}}(z_{i} - a_{i})f\omega = 0\) , whence

\[
\Theta_ {\boldsymbol {a}} (z _ {i} f) = \frac {n !}{(2 i \pi) ^ {n}} \int_ {\mathrm{U}} z _ {i} f \omega = \frac {n !}{(2 i \pi) ^ {n}} \int_ {\mathrm{U}} a _ {i} f \omega = a _ {i} \Theta_ {\boldsymbol {a}} (f)
\]

by formula (5). The result follows.

PROPOSITION 2. — Let \(\varrho_{1},\ldots,\varrho_{n}\) be real numbers \(>0\) and let \(\mathrm{U}\subset\mathbf{C}^{n}\) be the polydisc product of the open disks with center 0 and radius \(\varrho_{i}\) . Let

\[
\sum_ {(k _ {1}, \dots , k _ {n}) \in \mathbf {N} ^ {n}} c (k _ {1}, \dots , k _ {n}) \mathrm{X} _ {1} ^ {k _ {1}} \dots \mathrm{X} _ {n} ^ {k _ {n}} \in \mathrm{A} [ [ \mathrm{X} _ {1}, \dots , \mathrm{X} _ {n} ] ]
\]

a formal series with coefficients in A. Suppose that this series converges in U, and denote by f the holomorphic function in U which is its sum.

Let \(\boldsymbol{a}=(a_{1},\ldots,a_{n})\in\mathrm{A}^{n}\) such that \(\varrho(a_{i})<\varrho_{i}\) for \(1\leqslant i\leqslant n\) . Then \(\mathrm{Sp}^{n}(\boldsymbol{a})\subset\mathrm{U}\) , the family \((c(k_{1},\ldots,k_{n})a_{1}^{k_{1}}\cdots a_{n}^{k_{n}})\) of elements of A is absolutely summable, and

\[
\Theta_ {\boldsymbol {a}} (f) = \sum_ {(k _ {1}, \dots , k _ {n}) \in \mathbf {N} ^ {n}} c (k _ {1}, \dots , k _ {n}) a _ {1} ^ {k _ {1}} \dots a _ {n} ^ {k _ {n}}.
\]

For every character \(\chi\) of A and every integer \(i\) such that \(1 \leqslant i \leqslant n\) , one has \(|\chi(a_i)| \leqslant \varrho(a_i) < \varrho_i\) , hence \(\mathrm{Sp}^n(\boldsymbol{a}) \subset \mathrm{U}\) by definition of the simultaneous spectrum. Let \(z_1, \ldots, z_n\) be the restrictions to U of the coordinate functions on \(\mathbf{C}^n\) . Then the family \((c(k_1, \ldots, k_n)z_1^{k_1} \cdots z_n^{k_n})\) is summable in \(\mathcal{O}(\mathrm{U}; \mathrm{A})\) with sum \(f\) . Taking into account Lemma 10 and the continuity of the map \(\Theta_{\boldsymbol{a}}^{\mathrm{U}}\) , the family \((c(k_1, \ldots, k_n)a_1^{k_1} \ldots a_n^{k_n})\) is therefore summable in A with sum \(\Theta_{\boldsymbol{a}}(f)\) . For \(1 \leqslant i \leqslant n\) , let \(\lambda_i\) be a real number such that \(\varrho(a_i) < \lambda_i < \varrho_i\) . There exists \(\mathrm{M}_i < +\infty\) such that \(\|a_i^k\| \leqslant \mathrm{M}_i\lambda_i^k\) for every integer \(k \geqslant 0\) . We then have

\[
\sum_ {(k _ {1}, \dots , k _ {n}) \in \mathbf {N} ^ {n}} \| c (k _ {1}, \dots , k _ {n}) \| \| a _ {1} ^ {k _ {1}} \dots a _ {n} ^ {k _ {n}} \| \leqslant
\]

\[
\mathrm{M} _ {1} \dots \mathrm{M} _ {n} \sum_ {(k _ {1}, \dots , k _ {n}) \in \mathbf {N} ^ {n}} \| c (k _ {1}, \dots , k _ {n}) \| \lambda_ {1} ^ {k _ {1}} \dots \lambda_ {n} ^ {k _ {n}}
\]

which is finite by hypothesis, hence the family \((c(k_{1},\ldots,k_{n})a_{1}^{k_{1}}\cdots a_{n}^{k_{n}})\) is absolutely summable.

COROLLARY. --- Suppose A is nonzero. Let \(a \in C^{n} \subset A^{n}\) . We have \(\operatorname{Sp}_{\mathrm{A}}^{n}(\boldsymbol{a}) = \{\boldsymbol{a}\}\) . For every germ \(f \in \mathcal{O}(\{\boldsymbol{a}\}; \mathrm{A})\) , one has \(\Theta_{\boldsymbol{a}}(f) = f(\boldsymbol{a})\) .

PROPOSITION 3. --- Let B be a commutative unital Banach algebra and \(\varphi\) a continuous unital morphism from A to B. Let \(\boldsymbol{a} = (a_1, \ldots, a_n) \in \mathrm{A}^n\) . Set \(\boldsymbol{b} = (\varphi(a_1), \ldots, \varphi(a_n))\) , so that \(\mathrm{Sp}_{\mathrm{B}}^{n}(\boldsymbol{b}) \subset \mathrm{Sp}_{\mathrm{A}}^{n}(\boldsymbol{a})\) . For every \(f \in \mathcal{O}(\mathrm{Sp}_{\mathrm{A}}^{n}(\boldsymbol{a}); \mathrm{A})\) , one has

\[
\varphi (\Theta_ {\boldsymbol {a}} (f)) = \Theta_ {\boldsymbol {b}} (\varphi_ {*} (f)),
\]

where \(\varphi_{*}(f)\) denotes the germ of \(\varphi \circ f\) in a neighborhood of \(\mathrm{Sp}_{\mathrm{B}}^{n}(\boldsymbol{b})\) .

It suffices to prove that for every open neighborhood U of \(\mathrm{Sp}_{\mathrm{A}}^{n}(\boldsymbol{a})\) and every \(f \in \mathcal{O}(\mathrm{U}; \mathrm{A})\) , one has \(\varphi(\Theta_{\boldsymbol{a}}(f)) = \Theta_{\boldsymbol{b}}(\varphi \circ f)\) , where \(\varphi \circ f \in \mathcal{O}(\mathrm{U}; \mathrm{B})\) . Let \((h, u_{1}, \ldots, u_{n})\) be a sequence adapted to \(\boldsymbol{a}\), where the support of h is contained in U (Lemma 1 of I, p.~52 and Lemma 3 of I, p.~54). Denote by \(\omega\) the associated differential form. For every \(z \in C^{n}\) , one has

\[
\sum_ {j = 1} ^ {n} (z _ {j} - b _ {j}) \varphi (u _ {i} (\boldsymbol {z})) = \varphi \left(\sum_ {j = 1} ^ {n} (z _ {j} - a _ {j}) u _ {j} (\boldsymbol {z})\right) = 1 - h (\boldsymbol {z}),
\]

so that the sequence \((h,\varphi\circ u_{1},\ldots,\varphi\circ u_{n})\) is adapted to b. Let \(\omega'\) be the associated differential form. Denote by \(\mu\) the Lebesgue measure on \(C^{n}\) . Let \(f\in\mathcal{O}(U;A)\) . Write \(\psi\cdot\mu\) for the vector measure associated with the differential form \(f\omega\) . The vector measure associated with the differential form

\[
(\varphi \circ f) \omega^ {\prime} = (\varphi \circ f) d (\varphi \circ u _ {1}) \wedge d z _ {1} \wedge \dots \wedge d (\varphi \circ u _ {n}) \wedge d z _ {n}
\]

is equal to \((\varphi \circ \psi) \cdot \mu\) . Therefore, by formula (5), and INT, VI, §2, no. 2, prop. 2, we have

\[
\Theta_ {\boldsymbol {b}} (\varphi \circ f) = \frac {n !}{(2 i \pi) ^ {n}} \int_ {\mathrm{U}} (\varphi \circ f) \mu = \frac {n !}{(2 i \pi) ^ {n}} \varphi \left(\int_ {\mathrm{U}} \psi \mu\right) = \varphi (\Theta_ {\boldsymbol {a}} (f)),
\]

as was required.

COROLLARY 1. --- Let \(\chi \in \mathrm{X}(\mathrm{A})\) and \(\boldsymbol{a} \in \mathrm{A}^n\) . For every germ \(f \in \mathcal{O}(\mathrm{Sp}^n(\boldsymbol{a}))\) , one has \(\chi(\Theta_{\boldsymbol{a}}(f)) = f(\chi(a_1), \ldots, \chi(a_n))\) .

This is a consequence of Proposition 3, applied to the continuous unital morphism \(\chi\colon\mathbf{A}\to\mathbf{C}\) (th. 1 of I, p.~29), and of the corollary of prop. 2, applied to the Banach algebra C.

Remark. --- Suppose that the algebra A is semisimple. Let \(\boldsymbol{a} = (a_{1}, \ldots, a_{n}) \in \mathrm{A}^{n}\) . By prop. 8 of I, p.~38, the map \(\Theta_{\boldsymbol{a}}^{\mathrm{U}}\) is the unique map \(\varphi\) of \(\mathcal{O}(\mathrm{U})\) to A such that \(\chi(\varphi(f)) = f(\chi(a_{1}), \ldots, \chi(a_{n}))\) for all \(\chi \in \mathrm{X}(\mathrm{A})\) and every function \(f \in \mathcal{O}(\mathrm{U})\) .

COROLLARY 2. --- Let p be an integer ≥ 1. For every family \((f_{1},\ldots,f_{p})\) of elements of \(\mathcal{O}(\mathrm{Sp}^{n}(\boldsymbol{a}))\) , one has

\[
\operatorname{Sp} ^ {p} \left(\left(\Theta_ {\boldsymbol {a}} \left(f _ {1}\right), \dots , \Theta_ {\boldsymbol {a}} \left(f _ {p}\right)\right)\right) = \left(f _ {1}, \dots , f _ {p}\right) \left(\operatorname{Sp} ^ {n} (\boldsymbol {a})\right).
\]

In particular, for every \(f \in \mathcal{O}(\mathrm{Sp}^n(\boldsymbol{a}))\) , one has \(\mathrm{Sp}(\Theta_{\boldsymbol{a}}(f)) = f(\mathrm{Sp}^n(\boldsymbol{a}))\) .

This follows from Cor. 1 and the definition of the simultaneous spectrum.

Example. --- Let A be the complex Banach algebra of functions on the unit circle with absolutely convergent Fourier series (I, p.~19, example 8). Let \(\varphi \in\) A. Let \(f\) be a germ of a holomorphic function in a neighborhood of the set of values of \(\varphi\) . Then \(\psi = \Theta_{\varphi}(f)\) is an absolutely convergent Fourier series which for every \(u \in \mathbf{U}\) satisfies \(\psi(u) = f(\varphi(u))\) (cor. 1, applied to the characters \(\varphi \mapsto \varphi(u)\) ). In other words, the function \(f \circ \varphi\) on the unit circle also has an absolutely convergent Fourier series (“P. Lévy’s theorem”). This result generalizes Wiener’s theorem (I, p.~38, example 4), which concerns the case of the function \(f(z) = 1/z\) on \(\mathbf{C} - \{0\}\) when \(\varphi\) does not vanish.

\subsection{Approximation theorems}

In this issue, A is a commutative unital complex Banach algebra.

PROPOSITION 4. --- Let L be a compact polynomially convex subset of \(\mathbf{C}^{n}\) and let U be an open neighborhood of L. For every function \(f\in\mathcal{O}(U;A)\) , there exists a sequence of polynomial functions on \(\mathbf{C}^{n}\) with coefficients in A that converges to \(f|L\) into \(\mathcal{C}(L;A)\) .

We may assume that L is nonempty and that A is nonzero. Let P (resp. \(P_{0}\) ) be the set of restrictions to L of polynomial functions on \(C^{n}\) with coefficients in A (resp. with coefficients in C). Let B (resp. \(B_{0}\) ) be the Banach algebra closure of P (resp. of \(P_{0}\) )

into \(\mathcal{C}(\mathrm{L};\mathrm{A})\) . Denote by \(\iota\) the injection of A onto the normed subalgebra of B consisting of constant functions.

Let \(z_{1},\ldots,z_{n}\) be the restrictions to L of the coordinate functions on \(\mathbf{C}^{n}\) ; these are elements of \(\mathrm{B}_{0}\) , and by setting \(\boldsymbol{z}=(z_{1},\ldots,z_{n})\) it follows \(\mathrm{Sp}_{\mathrm{B}_{0}}^{n}(\boldsymbol{z})=\mathrm{L}\) according to prop. 15 of I, p.~47.

Let \(f \in \mathcal{O}(\mathrm{U};\mathrm{A})\) . By composition with \(\iota\) , the function \(f\) defines an element \(f_{\mathrm{B}} = \iota \circ f\) of \(\mathcal{O}(\mathrm{U};\mathrm{B})\) . Since \(\mathrm{Sp}_{\mathrm{B}}^{n}(\boldsymbol{z}) \subset \mathrm{Sp}_{\mathrm{B}_{0}}^{n}(\boldsymbol{z}) \subset \mathrm{U}\) , one can form the element \(b = \Theta_{\boldsymbol{z}}(f_{\mathrm{B}})\) of B. Let \(\boldsymbol{w} = (w_{1},\ldots ,w_{n}) \in \mathrm{L}\) , and let \(\varphi\) be the continuous unital morphism \(g \mapsto g(\boldsymbol{w})\) from B to A. We have \(\varphi \circ \iota = \mathrm{Id}_{\mathrm{A}}\) , so that \(\varphi \circ f_{\mathrm{B}} = f\) . Since \(\varphi (z_i) = w_i\) , prop. 3 of I, p.~66 implies \(\varphi (\Theta_z(f_{\mathrm{B}})) = \Theta_w(\varphi \circ f_{\mathrm{B}})\) . We therefore have

\[
b (\pmb {w}) = \varphi (b) = \varphi (\Theta_ {z} (f _ {\mathrm{B}})) = \Theta_ {\pmb {w}} (\varphi \circ f _ {\mathrm{B}}) = \Theta_ {\pmb {w}} (f) = f (\pmb {w}),
\]

by the corollary of prop. 2 of I, p.~65. Thus, we have \(f|_{\mathrm{L}} = b\) ; in particular, \(f|_{\mathrm{L}}\) belongs to B. This proves the proposition.

THEOREM 2 (Oka--Weil). --- Let K be a polynomially convex compact subset of \(C^{n}\) and P the set of germs in a neighborhood of K of polynomial functions on \(C^{n}\) with coefficients in A. Then P is dense in \(\mathcal{O}(K;A)\). More precisely, every element of \(\mathcal{O}(K;A)\) is the limit of a sequence of elements of P.

Consider an element of \(\mathcal{O}(\mathrm{K};\mathrm{A})\), the germ of a function \(f\in \mathcal{O}(\mathrm{U};\mathrm{A})\) , where U is an open neighborhood of K. By Lemma 7 of I, p.~48, there exists a compact neighborhood L of K contained in U which is polynomially convex. Let V be the interior of L; it is a neighborhood of K.

By the preceding proposition, there exists a sequence \((\mathrm{P}_{k})\) of polynomial functions on \(C^{n}\) with coefficients in A that converges to \(f|L\) into \(\mathcal{C}(L;A)\). In particular, the sequence \((\mathrm{P}_{k})\) converges to \(f|V\) into \(\mathcal{O}(V;A)\) .

By definition of the topology on \(\mathcal{O}(\mathrm{K};\mathrm{A})\) (cf. EVT, II, p. 29, prop. 5), the canonical mapping from \(\mathcal{O}(\mathrm{V};\mathrm{A})\) into \(\mathcal{O}(\mathrm{K};\mathrm{A})\) is continuous. Consequently, the sequence of germs in a neighborhood of K of the functions \(\mathrm{P}_k\) converges to the germ of \(f\) in the neighborhood of K in the space \(\mathcal{O}(\mathrm{K};\mathrm{A})\) , which is what had to be proved.

COROLLARY 1. --- Let U be an open neighborhood of K. Let \(u_{1}\) and \(u_{2}\) be continuous maps from \(\mathcal{O}(\mathrm{U};\mathrm{A})\) into a topological space X factoring through \(\mathcal{O}(\mathrm{K};\mathrm{A})\) . Then \(u_{1}=u_{2}\) if and only if \(u_{1}\) and \(u_{2}\) coincide on the set of restrictions to K of polynomial functions on \(\mathbf{C}^{n}\) with coefficients in A.

COROLLARY 2. --- Let E be a Banach space. Let K be a compact polynomially convex subset of \(\mathbf{C}^{n}\). Let P be the set of germs in a neighborhood of K of polynomial functions on \(\mathbf{C}^{n}\) with values in E. Then every element of \(\mathcal{O}(\mathrm{K};\mathrm{E})\) is the limit of a sequence of elements of P.

Equip E with the multiplication defined by ab = 0 for all a and b in E (example 1 of I, p.~17). This is a commutative Banach algebra. Let A be the unital commutative Banach algebra obtained from E by adjoining a unit element. Since the canonical map \(\mathcal{O}(\mathrm{K};\mathrm{A}) \to \mathcal{O}(\mathrm{K};\mathrm{E})\) is continuous, the assertion follows from the Oka-Weil theorem applied to the algebra A.

For n = 1 and A = C, we also have the following result, which will be made more precise by Corollary 2 of I, p.~150.

THEOREM 3 (Runge). --- Let K be a compact subset of C, and let Q be the set of germs of rational functions holomorphic in a neighborhood of K. Then Q is dense in \(\mathcal{O}(\mathrm{K})\) .

According to the definition of the topology on \(\mathcal{O}(\mathrm{K})\) , it suffices to show that for every open neighborhood U of K, and every compact subset L of U, every function \(f \in \mathcal{O}(\mathrm{U})\) is a limit of rational functions continuous on L. One may assume that L is a compact neighborhood of K.

Let Q' be the set of restrictions to L of rational functions on C that are continuous on L, and let C be the closure of Q' in \(\mathcal{C}(\mathrm{L})\) . This is an algebra without radical.

Let \(z \in \mathrm{C}\) be the identity map of L. Then C is the full closed subalgebra of \(\mathcal{C}(\mathrm{L})\) generated by \(z\) (lemma 2 of I, p.~6). We therefore have \(\mathrm{Sp}_{\mathrm{C}}(z) = \mathrm{Sp}_{\mathcal{C}(\mathrm{L})}(z) = \mathrm{L}\). We can then form the element \(c = \Theta_z(f)\) in C. Since C is radical-free, applying cor. 1 of I, p.~66 to the characters \(g \mapsto g(w)\) of C for all \(w \in \mathrm{L}\) shows that \(c\) coincides with the restriction of \(f\) to L. By definition of C, this proves that \(f|_{\mathrm{L}}\) is a uniform limit on L of elements of \(Q'\) , and this completes the proof of the theorem.

\subsection{Existence and uniqueness of the holomorphic functional calculus}

Suppose that A is a commutative unital complex Banach algebra.

DEFINITION 2. --- Let \(n \geqslant 1\) be an integer and \(\boldsymbol{a} \in \mathrm{A}^{n}\) . Let U be an open neighborhood of \(\mathrm{Sp}^{n}(\boldsymbol{a})\). We say that a family \(\boldsymbol{a}'\) is an envelope of \((\boldsymbol{a}, \mathrm{U})\) if \(\boldsymbol{a}' \in \mathbf{C}^{n+p}\) extends \(\boldsymbol{a}\) and if \(\mathrm{U} \times \mathbf{C}^{p}\) contains the polynomially convex envelope of \(\mathrm{Sp}^{n+p}(\boldsymbol{a}')\) .

Lemma 11. --- Let \(n \geqslant 1\) be an integer. Let \(\boldsymbol{a} \in \mathrm{A}^{n}\) . For every open neighborhood U of \(\mathrm{Sp}^{n}(\boldsymbol{a})\) , there exists an envelope of \((\boldsymbol{a}, \mathrm{U})\) .

Let \((a_{\lambda})_{\lambda\in\Lambda}\) be a family of elements of A extending the family a and topologically generating the unital Banach algebra A. Let \(\pi\) be the canonical projection from \(C^{\Lambda}\) onto \(C^{n}\) and let \(U'=\pi^{-1}(U)\) . Then \(U'\) is a neighborhood of \(\mathrm{Sp}^{\Lambda}((a_{\lambda}))\) , and \(\mathrm{Sp}^{\Lambda}((a_{\lambda}))\) is polynomially convex (I, p.~44, lemma 4). By lemma 6 of I, p.~47, there exists a finite subset \(\Lambda_{0}\) of \(\Lambda\) containing \(\{1,2,\ldots,n\}\) such that \(\mathrm{pr}_{\Lambda_{0}}(\mathrm{U}')\) contains the polynomially convex hull S of \(\mathrm{pr}_{\Lambda_{0}}(\mathrm{Sp}^{\Lambda}((a_{\lambda})_{\lambda\in\Lambda}))=\mathrm{Sp}^{\Lambda_{0}}((a_{\lambda})_{\lambda\in\Lambda_{0}})\) . Let \(p\geqslant0\) be the integer such that \(\Lambda_{0}\) has cardinality \(n+p\) and let j be a bijection from \(\{1,\ldots,n+p\}\) into \(\Lambda_{0}\) which coincides with the identity map on \(\{1,\ldots,n\}\). Since the projection of S is contained in U, the family \((a_{j(k)})_{1\leqslant k\leqslant n+p}\) is an envelope of \((\boldsymbol{a},\mathrm{U})\) .

PROPOSITION 5. --- The data of the mappings \(\Theta_{\boldsymbol{a}}\) , for \(n \geqslant 1\) and \(\boldsymbol{a} \in \mathrm{A}^{n}\) , is a holomorphic functional calculus on A, that is, conditions (CF1), (CF2), and (CF3) of I, p.~51 are satisfied.

Let \(n \geqslant 1\) be an integer and \(\boldsymbol{a} = (a_{1}, \ldots, a_{n}) \in \mathrm{A}^{n}\) . The map \(\Theta_{\boldsymbol{a}}\) satisfies \(\Theta_{\boldsymbol{a}}(z_{i}) = a_{i}\) for every i such that \(1 \leqslant i \leqslant n\) by lemma 10 of I, p.~64, which proves property (CF2). Lemma 9 of I, p.~63 implies property (CF3) for the maps \(\Theta_{\boldsymbol{a}}\) .

The map \(\Theta_{\pmb{a}}\) is A-linear and continuous (I, p.~61, no. 5). It satisfies \(\Theta_{\pmb{a}}(1) = 1\) (lemma 9 of I, p.~63). To verify condition (CF1), it remains to establish that \(\Theta_{\pmb{a}}\) is an algebra morphism. To do this, we will prove that \(\Theta_{\pmb{a}}^{\mathrm{U}}\) is an algebra morphism for every open neighborhood U of \(\mathrm{Sp}^{n}(\pmb{a})\) .

Suppose first that U contains the polynomially convex hull K of \(\mathrm{Sp}^{n}(\boldsymbol{a})\) . Let \(f_{1}\) and \(f_{2}\) be elements of \(\mathcal{O}(\mathrm{U};\mathrm{A})\) . There exists a sequence \((f_{1,k})\) (resp. \((f_{2,k})\)) of polynomial functions which converges to \(f_{1}\) (resp. to \(f_{2}\) ) in \(\mathcal{O}(\mathrm{K};\mathrm{A})\) (Theorem 2 of I, p.~68), hence in \(\mathcal{O}(\mathrm{Sp}^{n}(\boldsymbol{a});\mathrm{A})\) . For every integer \(k\) , one has

\[
\Theta_ {\boldsymbol {a}} ^ {\mathrm{U}} (f _ {1, k}) \Theta_ {\boldsymbol {a}} ^ {\mathrm{U}} (f _ {2, k}) = \Theta_ {\boldsymbol {a}} ^ {\mathrm{U}} (f _ {1, k} f _ {2, k})
\]

by Lemma 10 of I, p.~64, whence \(\Theta_{a}^{\mathrm{U}}(f_{1})\Theta_{a}^{\mathrm{U}}(f_{2}) = \Theta_{a}^{\mathrm{U}}(f_{1}f_{2})\) by passing to the limit.

Consider the general case. Let \(a' \in C^{n+p}\) be an envelope of \((\boldsymbol{a}, U)\) (lemma 11) and \(\pi: U \times C^{p} \to U\) the canonical projection. Since \(U \times C^{p}\) contains the polynomially convex envelope of \(\mathrm{Sp}^{n+p}(\boldsymbol{a}')\) , one has

\[
\Theta_ {\boldsymbol {a} ^ {\prime}} ^ {\mathrm{U} \times \mathbf {C} ^ {p}} (f _ {1} \circ \pi) \Theta_ {\boldsymbol {a} ^ {\prime}} ^ {\mathrm{U} \times \mathbf {C} ^ {p}} (f _ {2} \circ \pi) = \Theta_ {\boldsymbol {a} ^ {\prime}} ^ {\mathrm{U} \times \mathbf {C} ^ {p}} (f _ {1} f _ {2} \circ \pi)
\]

for \(f_1\) and \(f_2\) in \(\mathcal{O}(\mathrm{U};\mathrm{A})\) by the first case. As, for every function \(f \in \mathcal{O}(\mathrm{U};\mathrm{A})\) , one has \(\Theta_{\boldsymbol{a}'}^{\mathrm{U}\times \mathbf{C}^p}(f\circ \pi) = \Theta_{\boldsymbol{a}}^{\mathrm{U}}(f)\) (condition (CF3) previously proved), the conclusion follows, and hence condition (CF1).

We can now prove Theorem 1 of I, p.~51. Prop. 5 shows that the family of maps \((\Theta_{a})_{a}\) is a holomorphic functional calculus on A. It therefore remains only to establish the uniqueness of the holomorphic functional calculus on A.

Let \((\Psi_{a})_{a}\) be a family of maps defined for every integer \(n \geqslant 1\) and every \(a \in A^{n}\) and satisfying conditions (CF1), (CF2), (CF3) of the holomorphic functional calculus on A (I, p.~51). It suffices to prove that for every integer \(n \geqslant 1\), every \(a \in A^{n}\), and every open neighborhood U of a, we have \(\Theta_{a}^{U} = \Psi_{a}^{U}\) .

Let \(n \geqslant 1\) and \(\boldsymbol{a} = (a_{1}, \ldots, a_{n}) \in \mathrm{A}^{n}\) . Let U be an open neighborhood of \(\mathrm{Sp}^{n}(\boldsymbol{a})\) . Suppose first that U contains the polynomially convex hull K of \(\mathrm{Sp}^{n}(\boldsymbol{a})\) . The morphisms \(\Theta_{a}^{U}\) and \(\Psi_{a}^{U}\) coincide on polynomial functions by properties (CF1) and (CF2). By the corollary of Theorem 2 of I, p.~68 and the continuity property (CF1), these morphisms are therefore equal.

Let us prove the general case. Let \(a' \in C^{n+p}\) be an envelope of \((\boldsymbol{a}, U)\) and \(\pi: U \times C^{p} \to U\) the canonical projection. We have

\[
\Theta_ {\boldsymbol {a}} ^ {\mathrm{U}} = \Theta_ {\boldsymbol {a} ^ {\prime}} ^ {\mathrm{U} \times \mathbf {C} ^ {p}} \circ \pi^ {*} = \Psi_ {\boldsymbol {a} ^ {\prime}} ^ {\mathrm{U} \times \mathbf {C} ^ {p}} \circ \pi^ {*} = \Psi_ {\boldsymbol {a}} ^ {\mathrm{U}}
\]

by property (CF3) and the preceding case. This concludes the proof of Theorem 1 of I, p.~51.

Note that Th. 2 of I, p.~68 also entails the following uniqueness result:

PROPOSITION 6. --- Let \(\boldsymbol{a} \in \mathrm{A}^n\) . We suppose \(\mathrm{Sp}^n(\boldsymbol{a})\) polynomially convex. Let \(z_1, \ldots, z_n\) be the germs in a neighbourhood of \(\mathrm{Sp}^n(\boldsymbol{a})\) of the coordinate functions on \(\mathbf{C}^n\). Then the map \(\Theta_{\boldsymbol{a}}\) is the unique continuous morphism of unital algebras \(\varphi\) of \(\mathcal{O}(\mathrm{Sp}^n(\boldsymbol{a}); \mathrm{A})\) to A such that \(\varphi(z_1) = a_1, \ldots, \varphi(z_n) = a_n\) .

Lemma 10 of I, p.~64 and the corollary of Proposition 2 of I, p.~65 justify the following notations for the holomorphic functional calculus. Let \(n \geqslant 1\) and \(\boldsymbol{a} \in \mathrm{A}^n\) . For every germ \(f \in \mathcal{O}(\mathrm{K};\mathrm{A})\) (resp. for every holomorphic function \(f \in \mathcal{O}(\mathrm{U};\mathrm{A})\) on an open neighbourhood U of \(\mathrm{Sp}^n (\boldsymbol{a})\) ), we set

\[
f (\boldsymbol {a}) = \Theta_ {\boldsymbol {a}} (f).\tag{6}
\]

This notation is consistent with the notation introduced in A, IV, p.~4, n°3, if f is a polynomial, by properties (CF1) and (CF2). Properties (CF2) and (CF3) of I, p.~51 can then be written

\[
z _ {i} (\boldsymbol {a}) = a _ {i}, \quad 1 \leqslant i \leqslant n, \quad (f \circ \pi_ {m, n}) (\boldsymbol {a}) = f (\pi_ {m, n} (\boldsymbol {a})).
\]

\subsection{Substitution in the functional calculus}

With the notations introduced above, the statements of Cor. 1 of I, p.~66 and Cor. 2 of I, p.~67 become respectively

\[
\chi (g (\boldsymbol {a})) = g \left(\chi \left(a _ {1}\right), \dots , \chi \left(a _ {n}\right)\right), \quad \operatorname{Sp} (g (\boldsymbol {a})) = g \left(\operatorname{Sp} ^ {n} (\boldsymbol {a})\right)
\]

for \(f \in \mathcal{O}(\mathrm{Sp}^n(\boldsymbol{a});\mathrm{A})\) , \(\chi \in \mathrm{X}(\mathrm{A})\) and \(g \in \mathcal{O}(\mathrm{Sp}^n(\boldsymbol{a}))\) .

We shall now prove a more general substitution property.

THEOREM 4. --- Let A be a commutative unital complex Banach algebra, let \(n \geqslant 1\) be an integer and \(\boldsymbol{a} = (a_1, \ldots, a_n) \in \mathrm{A}^n\) . Let \(\boldsymbol{f} = (f_1, \ldots, f_p)\) where \(f_1, \ldots, f_p\) are elements of \(\mathcal{O}(\mathrm{Sp}^n(\boldsymbol{a}))\) . The image of \(\mathrm{Sp}^n(\boldsymbol{a})\) by the map \(\boldsymbol{z} \mapsto \boldsymbol{f}(\boldsymbol{z}) = (f_1(\boldsymbol{z}), \ldots, f_p(\boldsymbol{z}))\) is equal to \(\mathrm{Sp}^p(\boldsymbol{f}(\boldsymbol{a}))\) .

For every \(g \in \mathcal{O}(\mathrm{Sp}^{p}(\boldsymbol{f}(\boldsymbol{a}));\mathrm{A})\) , the composed germ \(g \circ f\) is an element of \(\mathcal{O}(\mathrm{Sp}^{n}(\boldsymbol{a});\mathrm{A})\) and we have \(g(\boldsymbol{f}(\boldsymbol{a})) = (g \circ \boldsymbol{f})(\boldsymbol{a})\) .

The first assertion concerning the image of \(\mathrm{Sp}^{n}(\boldsymbol{a})\) follows from Cor. 2 of I, p.~67. To prove the second, we shall use the following lemma.

Lemma 12. --- Let K be the polynomially convex hull of \(\mathrm{Sp}^p (\pmb {f}(\pmb {a}))\) . We have \(g(\pmb {f}(\pmb {a})) = (g\circ \pmb {f})(\pmb {a})\) for every germ \(g\in \mathcal{O}(\mathrm{K};\mathrm{A})\) .

Let \(\Psi\) be the map from \(\mathcal{O}(\mathrm{K};\mathrm{A})\) to A such that \(\Psi(g) = (g\circ f)(\boldsymbol{a})\) . It is a continuous unital morphism, such that \(\Psi(z_j) = f_j(\boldsymbol{a})\) , where \(z_{j}\) is the germ of the \(j\) -th coordinate function on \(\mathbf{C}^p\). When \(g\) is the germ of a polynomial function, we therefore have \(\Psi(g) = g(\boldsymbol{f}(\boldsymbol{a}))\). By th. 2 of I, p.~68, this formula remains valid for all \(g\in \mathcal{O}(\mathrm{K};\mathrm{A})\) .

Now let us prove the theorem. Let V be an open neighborhood of \(\mathrm{Sp}^{p}(\boldsymbol{f}(\boldsymbol{a}))\) and \(\widetilde{g}\in\mathcal{O}(V;A)\) be a holomorphic function whose germ in a neighborhood of \(\mathrm{Sp}^{p}(\boldsymbol{f}(\boldsymbol{a}))\) is equal to g. Let \(b\in C^{p+q}\) be an envelope of \((\boldsymbol{f}(\boldsymbol{a}),\mathrm{V})\) (Lemma 11 of I, p.~70) and \(\pi\colon V\times C^{q}\to V\) the canonical projection.

Let \(\widetilde{f}_1, \ldots, \widetilde{f}_p\) be holomorphic functions whose germs are \(f_1, \ldots, f_p\) and let U be an open neighborhood of \(\mathrm{Sp}^n(\boldsymbol{a})\) such that \((\widetilde{f}_1, \ldots, \widetilde{f}_p)(\mathrm{U}) \subset \mathrm{V}\) . Let \(\pi'\) be the canonical projection from \(\mathrm{U} \times \mathbf{C}^q\) onto U. Denote by \(z_{n+1}, \ldots, z_{n+q}\) the \(q\) last coordinate functions on \(\mathbf{C}^{n+q}\) . Denote by \(h = \widetilde{g} \circ (\widetilde{f}_1, \ldots, \widetilde{f}_p)\) and

\[
\boldsymbol {c} = (a _ {1}, \dots , a _ {n}, b _ {p + 1}, \dots , b _ {p + q}) \in \mathrm{A} ^ {n + q}.
\]

The map \(\widetilde{g} \circ \pi\) is holomorphic in the open neighborhood \(V \times C^q\) of the polynomially convex hull L of \(Sp^{p+q}(b)\). By Lemma 12, applied to \(c\) and to the germs in the neighborhood of L of the functions

\[
(\widetilde {f} _ {1} \circ \pi^ {\prime}, \ldots , \widetilde {f} _ {p} \circ \pi^ {\prime}, z _ {n + 1}, \ldots , z _ {n + q}),
\]

and to the germ of \(\widetilde{g}\circ\pi\) , one has

\[
(g \circ \pi) \left((f _ {1} \circ \pi^ {\prime}) (\boldsymbol {c}), \dots , (f _ {p} \circ \pi^ {\prime}) (\boldsymbol {c}), z _ {n + 1} (\boldsymbol {c}), \dots , z _ {n + q} (\boldsymbol {c})\right) = (h \circ \pi^ {\prime}) (\boldsymbol {c}).
\]

Since \(\pi'(\boldsymbol{c}) = \boldsymbol{a}\) , one has \((h \circ \pi')(\boldsymbol{c}) = h(\boldsymbol{a})\) and \((f_i \circ \pi')(\boldsymbol{c}) = f_i(\boldsymbol{a})\) for \(1 \leqslant i \leqslant p\) (property (CF3) of the holomorphic functional calculus). Moreover, since \(z_{n+j}(\boldsymbol{c}) = b_{p+j}\) for \(1 \leqslant j \leqslant q\). By property (CF2), we have

\[
(g \circ \pi) (f _ {1} (\boldsymbol {a}), \dots , f _ {p} (\boldsymbol {a}), b _ {p + 1}, \dots , b _ {p + q}) = h (\boldsymbol {a}),
\]

from which we deduce \(g(f_{1}(\boldsymbol{a}), \ldots, f_{p}(\boldsymbol{a})) = h(\boldsymbol{a})\) by applying property (CF3) again.

\subsection{Holomorphic functional calculus in one variable}

THEOREM 5. --- Let A be a unital Banach algebra, not necessarily commutative. Let a be an element of A and z the germ of the identity function of C in a neighborhood of SpA(a). There exists a unique unital continuous morphism \(\varphi_{a}\) of \(\mathcal{O}(\mathrm{Sp}_{\mathrm{A}}(a))\) to A such that \(\varphi_{a}(z)=a\) .

The image of \(\varphi_{a}\) is contained in the full closed subalgebra of A generated by \(a\) . In particular, it is contained in the bicommutant of \(a\) .

Let us prove the existence of the morphism \(\varphi_{a}\) . Let B be the closed full subalgebra of A generated by \(a\) . It is commutative, and one has \(\mathrm{Sp}_{\mathrm{B}}(a) = \mathrm{Sp}_{\mathrm{A}}(a)\) (I, p.~5, no. 5). The mapping \(\Theta_{a}\) of holomorphic functional calculus on B is a continuous unital morphism from \(\mathcal{O}(\mathrm{Sp}_{\mathrm{B}}(a))\) to B such that \(\Theta_{a}(z) = a\) (Theorem 1 of I, p.~51). The composed morphism of \(\Theta_{a}\) and of the canonical injection of B into A is a continuous unital morphism \(\varphi_{a}\) of \(\mathcal{O}(\mathrm{Sp}_{\mathrm{A}}(a))\) into A such that the image of \(z\) is \(a\) .

Let us prove uniqueness. Let \(\varphi_{a}^{\prime}\) be a continuous unital morphism from \(\mathcal{O}(\mathrm{Sp}_{\mathrm{A}}(a))\) to A such that \(\varphi_{a}^{\prime}(z)=a\) . Then \(\varphi_{a}\) and \(\varphi_{a}^{\prime}\) coincide on the set of germs of polynomials in a neighborhood of \(\mathrm{Sp}_{\mathrm{A}}(a)\) , hence on the set of germs of holomorphic rational fractions in a neighborhood of \(\mathrm{Sp}_{\mathrm{A}}(a)\) . Now these germs are dense in \(\mathcal{O}(\mathrm{Sp}_{\mathrm{A}}(a))\) (I, p.~69, Th. 3). This implies that \(\varphi_{a}=\varphi_{a}^{\prime}\) .

The construction of \(\varphi_{a}\) shows that its image is contained in the commutative subalgebra B, which is contained in the bicommutant of a (I, p.~6).

If the radical of the algebra A is zero, the uniqueness of the morphism \(\varphi_{a}\) is valid without requiring it to be continuous (cf.~prop. 9 of I, p.~40). This is not the case in general, cf.~G. R. ALLAN, Embedding the algebra of formal power series in a Banach algebra, Proc. London Math. Soc. (3) 25 (1972), 329--340.

For every Banach algebra A, every element \(a\) of A and every germ \(f \in \mathcal{O}(\mathrm{Sp}_{\mathrm{A}}(a))\) , we denote by \(f(a)\) the element \(\varphi_{a}(f)\) of Theorem 5. If A is a commutative Banach algebra, this element \(f(a)\) coincides with the element \(f(a)\) provided by the holomorphic functional calculus on a commutative Banach algebra (Theorem 1 of I, p.~51).

Let B be the closed full subalgebra of A generated by \(a\) , so that \(\mathrm{Sp}_{\mathrm{A}}(a) = \mathrm{Sp}_{\mathrm{B}}(a)\) . The element \(f(a)\) of A belongs to B, and coincides with the element \(f(a)\) computed relative to the algebra B.

PROPOSITION 7. --- Let A and B be unital Banach algebras and \(\varphi\) a continuous unital morphism from A to B. Let \(a \in A\) . Then \(\mathrm{Sp}_{\mathrm{B}}(\varphi(a)) \subset \mathrm{Sp}_{\mathrm{A}}(a)\) and we have \(\varphi(f(a)) = f(\varphi(a))\) for all \(f \in \mathcal{O}(\mathrm{Sp}_{\mathrm{A}}(a))\) . In particular, for every \(\chi \in \mathsf{X}(\mathrm{A})\) , one has \(\chi(f(a)) = f(\chi(a))\) .

This follows from prop. 3 of I, p.~66.

PROPOSITION 8. --- Let A be a unital Banach algebra and \(a \in A\) . Let \(f \in \mathcal{O}(\mathrm{Sp}(a))\) . We have \(f(\mathrm{Sp}_{\mathrm{A}}(a)) = \mathrm{Sp}_{\mathrm{A}}(f(a))\) . Moreover, for every \(g \in \mathcal{O}(\mathrm{Sp}_{\mathrm{A}}(f(a)))\) , one has \(g \circ f \in \mathcal{O}(\mathrm{Sp}_{\mathrm{A}}(a))\) and \(g(f(a)) = (g \circ f)(a)\) .

This follows from th. 4.

PROPOSITION 9. --- Let A be a unital Banach algebra and \(a \in A\) . Let U be an open neighborhood of \(\mathrm{Sp}_{\mathrm{A}}(a)\) and \(f \in \mathcal{O}(\mathrm{U})\) . Let moreover V be a compact neighborhood of \(\mathrm{Sp}_{\mathrm{A}}(a)\) contained in U such that V is a piece of U with oriented boundary \(\partial\mathrm{V}\) .

For every integer \(n \geqslant 0\) , the map \(z \mapsto f(z)(z - a)^{-n-1}\) is continuous on \(\partial V\); the differential form \(z \mapsto f(z)(z - a)^{-n-1} dz\) is integrable on \(\partial V\), and we have

\[
f ^ {(n)} (a) = \frac {n !}{2 i \pi} \int_ {\partial \mathrm{V}} f (z) (z - a) ^ {- n - 1} d z\tag{7}
\]

where \(f^{(n)}\in\mathcal{O}(\mathrm{U})\) is the n-th derivative of \(f\) .

Proceed by induction on \(n\). When \(n = 0\) , the result follows from prop. 1 of I, p.~61. Now suppose that the assertion of the proposition is true for the integer \(n \geqslant 0\) . Let \(g \in \mathcal{O}(\mathbf{C} - \mathrm{Sp}_{\mathrm{A}}(a); \mathrm{A})\) be the holomorphic function defined by \(g(z) = (z - a)^{-n-1} f(z)\). The differential form \(g'(z) dz = dg\) is of class \(\mathrm{C}^1\) ; since the piece V is compact, Stokes' formula (VAR, R2, p.~47, 11.2.3) implies

\[
\int_ {\partial \mathrm{V}} g ^ {\prime} (z) d z = \int_ {\partial \mathrm{V}} d g = 0.
\]

Since \(g'(z) = (z - a)^{-n-1} f'(z) - (n+1)(z - a)^{-n-2} f(z)\) , from which one deduces

\[
\int_ {\partial \mathrm{V}} f ^ {\prime} (z) (z - a) ^ {- n - 1} d z = (n + 1) \int_ {\partial \mathrm{V}} f (z) (z - a) ^ {- n - 2} d z.
\]

Applying the induction hypothesis to \(f'\) , one therefore obtains

\[
\frac {2 i \pi}{n !} f ^ {(n + 1)} (a) = (n + 1) \int_ {\partial \mathrm{V}} f (z) (z - a) ^ {- n - 2} d z,
\]

which is the assertion of the proposition for the integer \(n+1\) . This concludes the proof.

PROPOSITION 10. --- Let A be a unital Banach algebra and U an open subset of C.

\begin{enumerate}
\def\labelenumi{\alph{enumi})}
\tightlist
\item
  The set \(\Omega\) of \(a \in \mathrm{A}\) such that \(\operatorname{Sp}_{\mathrm{A}}(a) \subset \mathrm{U}\) is open in A;
\item
  Let \(f \in \mathcal{O}(\mathrm{U})\) . The map \(a \mapsto f(a)\) of \(\Omega\) in A is holomorphic, and in particular continuous.
\end{enumerate}

Let \(a \in \Omega\) . There exists a compact neighborhood V of \(\mathrm{Sp}_{\mathrm{A}}(a)\) contained in U which is a piece of U (VAR, R2, p.~46 and p.~47, 11.1.3, d)).

Since the resolvent of \(a\) tends to 0 at infinity (th. 1, \(c\) ) of I, p.~24), the map \(z \mapsto \| (z - a)^{-1}\|\) is bounded on \(\mathbf{C} - \mathring{\mathrm{V}}\) . Let M denote its upper bound. If \(h \in \mathrm{A}\) is such that \(\| h\| \leqslant (2\mathrm{M})^{-1}\) and if \(z \in \mathbf{C} - \mathring{\mathrm{V}}\) , one has

\[
z - (a + h) = (1 - h (z - a) ^ {- 1}) (z - a)
\]

and \(\|h(z-a)^{-1}\|\leqslant\frac{1}{2}\) , hence \(z-(a+h)\) is invertible and its inverse satisfies

\[
(z - (a + h)) ^ {- 1} = (z - a) ^ {- 1} \sum_ {n = 0} ^ {\infty} (h (z - a) ^ {- 1}) ^ {n}\tag{8}
\]

with \(\| (h(z - a)^{-1})^n\| \leqslant 2^{-n}\) (prop. 2 of I, p.~22). Thus, \(\mathrm{Sp}_{\mathrm{A}}(a + h)\) is contained in V, hence in U, which proves that \(\Omega\) is open in A.

Let \(f \in \mathcal{O}(\mathrm{U})\) . Denote by \(m\) the supremum of \(|f(z)|\) for \(z \in \partial \mathrm{V}\) . Let \(a \in \mathrm{A}\) . For every \(h \in \mathrm{A}\) such that \(\| h \| \leqslant (2\mathrm{M})^{-1}\) , one has

\[
f (a + h) = \frac {1}{2 i \pi} \int_ {\partial \mathrm{V}} f (z) (z - (a + h)) ^ {- 1} d z
\]

(prop. 9). The series (8) converges uniformly on the boundary of V, hence

\[
f (a + h) = \sum_ {n = 0} ^ {+ \infty} f _ {a, n} (h)
\]

where the map \(f_{a,n}\) from A to A is defined by

\[
f _ {a, n} (h) = \frac {1}{2 i \pi} \int_ {\partial \mathrm{V}} f (z) (z - a) ^ {- 1} (h (z - a) ^ {- 1}) ^ {n} d z.
\]

For every \(n \in N\) , the function \(f_{a,n}\) is a continuous homogeneous polynomial function of degree n. Moreover, it follows that

\[
\left\| f _ {a, n} (h) \right\| \leqslant \frac {m M}{\pi} \left(\int_ {\partial V} \| d z \|\right) 2 ^ {- (n + 1)}
\]

(INT, VI, §2, no. 3, prop. 5). The series \(\sum_{n} f_{a,n}(h)\) is therefore absolutely convergent for \(\|h\| \leqslant (2M)^{-1}\). This proves that the map which to \(a\) , associates \(f(a)\) is holomorphic on \(\Omega\) (VAR, R1, p.~26, 3.2.1).

PROPOSITION 11. --- Let A be a unital Banach algebra, \(a \in A\) and U an open neighborhood of \(\mathrm{Sp}_{\mathrm{A}}(a)\) . Denote by \(\delta\) the distance from \(\mathrm{Sp}_{\mathrm{A}}(a)\) to \(\mathbf{C} - \mathbf{U}\) . Let \(f \in \mathcal{O}(\mathbf{U})\) .

\begin{enumerate}
\def\labelenumi{\alph{enumi})}
\item
  For every real number η such that 0 \textless{} η \textless{} δ, there exists a real number C ≥ 0 such that \(\|f^{(n)}(a)\| \leqslant \mathrm{Cn}!\eta^{-n}\) for every integer \(n \in N\) ;
\item
  If \(b \in A\) commutes with \(a\) and if \(\varrho(b) < \delta\) , one has \(\mathrm{Sp}_{\mathrm{A}}(a + b) \subset \mathrm{U}\) , and
\end{enumerate}

\[
f (a + b) = \sum_ {n = 0} ^ {\infty} \frac {f ^ {(n)} (a)}{n !} b ^ {n},
\]

where the series converges absolutely.

Let \(\eta\) be a real number such that \(0 < \eta < \delta\) . Denote by \(\varepsilon = \delta - \eta > 0\) . Let K be the compact neighborhood of \(\mathrm{Sp}_{\mathrm{A}}(a)\) consisting of the points of C whose distance to \(\mathrm{Sp}_{\mathrm{A}}(a)\) is \(\leqslant \varepsilon / 2\) . Since \(f\) is holomorphic in every open disk of radius \(\eta + \varepsilon / 2\) whose center belongs to K, there exists, by the Cauchy inequalities (VAR, R1, p.~29, 3.3.4), a real number \(C \geqslant 0\) such that

\[
\sup _ {z \in K} \frac {| f ^ {(n)} (z) |}{n !} \leqslant \frac {C}{\eta^ {n}}
\]

for every integer \(n \geqslant 0\) . Then assertion a) follows from prop. 1 of I, p.~61 applied to \(f^{(n)}\) and to a piece V contained in K.

Let \(b\) be an element of A commuting with \(a\) such that \(\varrho(b) < \delta\) . By replacing A by the closed full subalgebra B generated by \(a\) and \(b\) , which satisfies \(\mathrm{Sp}_{\mathrm{A}}(a) = \mathrm{Sp}_{\mathrm{B}}(a)\) and \(\mathrm{Sp}_{\mathrm{A}}(a + b) = \mathrm{Sp}_{\mathrm{B}}(a + b)\) , one reduces, in order to prove \(b\) ), to the case where A is commutative.

Since \(\varrho(b) < \delta\) , one can choose \(\eta\) such that \(\varrho(b) < \eta < \delta\) . Let \(V_1\) be the set of points of \(\mathbf{C}\) whose distance to \(\mathrm{Sp}_{\mathrm{A}}(a)\) is \(< \delta - \eta\) , and \(V_2\) be the open disk with center 0 and radius \(\eta\) in \(\mathbf{C}\) . Let \(g\) be the map \((z_1, z_2) \mapsto z_1 + z_2\) from \(V_1 \times V_2\) into U. Then \(h = f \circ g\) is the map \((z_1, z_2) \mapsto f(z_1 + z_2)\) from \(V_1 \times V_2\) into \(\mathbf{C}\) . We have \(\mathrm{Sp}_{\mathrm{A}}^{2}(a, b) \subset V_1 \times V_2\) , hence \(\mathrm{Sp}_{\mathrm{A}}(a + b) \subset U\) (cf. Cor. 2 of I, p. 67), and moreover \(f(a + b) = h(a, b)\) . By th. 4 of I, p.~72. Now, in the space \(\mathcal{O}(V_1 \times V_2)\) , one has

\[
h (z _ {1}, z _ {2}) = \sum_ {n \geqslant 0} \frac {f ^ {(n)} (z _ {1})}{n !} z _ {2} ^ {n},
\]

therefore the series

\[
\sum_ {n \geqslant 0} \frac {f ^ {(n)} (a)}{n !} b ^ {n}
\]

converges in A and its sum is \(h(a,b)=f(a+b)\). Moreover, this series is absolutely convergent by assertion (a).

\subsection{Exponential and logarithm}

We denote by exp the complex exponential function from C to C (FVR, III, p.~8, def. 2). It is differentiable and satisfies \(\exp' = \exp\) (FVR, III, p.~9, (26)); therefore it is holomorphic in C. Let A be a unital Banach algebra and \(a\) an element of A. By prop. 2 of I, p.~65 and formula (9) of FVR, III, p.~16, we have

\[
\exp (a) = \sum_ {n = 0} ^ {\infty} \frac {a ^ {n}}{n !}.\tag{9}
\]

Since \(\|a^{n}\|\leqslant\|a\|^{n}\) we see that \(\|\exp(a)\|\leqslant\exp(\|a\|)\) and that the series (9) converges uniformly in every ball of A. The map \(a\mapsto\exp(a)\) from A to A is holomorphic (prop. 10 of I, p.~76). One also sometimes denotes \(e^{a}\) the exponential of \(a\in A\) .

When \(a\) is an endomorphism of a Banach space E, the exponential \(\exp(a)\) thus defined in the Banach algebra \(\mathcal{L}(\mathrm{E})\) coincides with that defined in FVR, IV, p.~27, def. 1, according to loc. cit. prop. 7 (3).

For every element b of A that is permutable with a, we also have

\[
\exp (a + b) = \sum_ {n = 0} ^ {\infty} \frac {b ^ {n}}{n !} \exp (a),
\]

(prop. 11 of I, p.~77), whence

\[
\exp (a + b) = \exp (a) \cdot \exp (b).\tag{10}
\]

In particular, \(\exp(a)\) is invertible and

\[
\exp (a) ^ {- 1} = \exp (- a).\tag{11}
\]

Let B be the set of \(z \in C\) such that \(-\pi < Jz < \pi\) . Let F be the complement in C of the interval \(R_{-}\). The restriction of the exponential to B induces, by passing to subspaces, a bijection from B onto F (FVR, III, p.~10, n°7), whose inverse bijection will be denoted by log.

If \(a \in A\) is such that \(\mathrm{Sp}_{\mathrm{A}}(a) \subset \mathrm{F}\) , one can form the element \(\log(a)\) of A. We have \(\mathrm{Sp}_{\mathrm{A}}(\log(a)) \subset \mathrm{B}\) , and

\[
\exp (\log (a)) = a\tag{12}
\]

by prop. 8 of I, p.~75. Conversely, let \(b\) be an element of A such that \(\mathrm{Sp}_{\mathrm{A}}(b) \subset \mathrm{B}\) . We have \(\mathrm{Sp}_{\mathrm{A}}(\exp(b)) \subset \mathrm{F}\) and

\[
\log (\exp (b)) = b\tag{13}
\]

(loc. cit.).

In particular, if \(a \in A\) is such that \(\varrho(a) < 1\) , one has \(\mathrm{Sp}_{\mathrm{A}}(1 - a) \subset \mathrm{F}\) and we can form \(\log(1 - a)\) . For \(n \geqslant 1\) , the \(n\) th derivative of \(z \mapsto \log(1 - z)\) is \(z \mapsto -(n - 1)! (1 - z)^{-n}\). The power series expansion of \(z \mapsto \log(1 - z)\) at the point 0 is therefore

\[
\log (1 - z) = - \sum_ {n = 1} ^ {\infty} \frac {z ^ {n}}{n},
\]

valid for \(|z| < 1\) (VAR, R1, p.~30, 3.3.9). By prop. 2 of I, p.~65, we obtain

\[
\log (1 - a) = - \sum_ {n = 1} ^ {\infty} \frac {a ^ {n}}{n}.\tag{14}
\]

PROPOSITION 12. --- Let A be a commutative unital Banach algebra. The image of the exponential map is the connected component of the identity in the group G of invertible elements of A.

Formulas (10) and (11) prove that \(\exp(A)\) is a subgroup of G. By the preceding (see formula (12)), this subgroup contains the open ball of center 1 and radius 1. It is therefore an open subgroup, and consequently closed, of G. Moreover, A is connected and the map \(a \mapsto \exp(a)\) is continuous, so that \(\exp(A)\) is connected. Therefore \(\exp(A)\) is the identity component of G.

\subsection{Partitions of the character space}

PROPOSITION 13. --- Let A be a commutative unital Banach algebra. Let U₁ and U₂ be open subsets of X(A) forming a partition of X(A). Then there exists a unique idempotent j in A such that the Gelfand transform G(j) equals 1 on U₁ and 0 on U₂.

Let us identify the space \(\mathsf{X}(\mathsf{A})\) with a compact subset of \(\mathbf{C}^{\mathrm{A}}\) by the map \(\chi \mapsto (\chi (a))_{a\in \mathrm{A}}\) (cf. no. 6 of I, p. 6 and cor. of th. 1 of I, p. 29). The subsets \(\mathrm{U}_1\) and \(\mathrm{U}_2\) of the uniform space \(\mathbf{C}^{\mathrm{A}}\) are compact and disjoint.

By TG, II, p.~31, prop. 4, there exists a finite subset M of A and disjoint open sets \(\mathrm{V}_{1}\) and \(\mathrm{V}_{2}\) of \(\mathbf{C}^{\mathrm{M}}\) such that

\[
p (\mathrm{U} _ {1}) \subset \mathrm{V} _ {1}, \qquad p (\mathrm{U} _ {2}) \subset \mathrm{V} _ {2},
\]

where \(p\) is the canonical projection of \(\mathbf{C}^{\mathrm{A}}\) on \(\mathbf{C}^{\mathrm{M}}\) .

Let \(a_1, \ldots, a_n\) be the distinct elements of M, and identify \(\mathbf{C}^{\mathrm{M}}\) with \(\mathbf{C}^n\) . We have \(\mathrm{Sp}_{\mathrm{A}}^{n}(a_{1}, \ldots, a_{n}) \subset p(\mathsf{X}(\mathrm{A})) \subset \mathrm{V}_{1} \cup \mathrm{V}_{2}\) since \(\mathrm{U}_{1} \cup \mathrm{U}_{2} = \mathsf{X}(\mathrm{A})\) . Let \(f\) be the function on \(\mathrm{V}_{1} \cup \mathrm{V}_{2}\) equal to 1 on \(\mathrm{V}_{1}\) and to 0 on \(\mathrm{V}_{2}\) . We have \(f \in \mathcal{O}(\mathrm{V}_{1} \cup \mathrm{V}_{2})\) . Set \(j = f(a_{1}, \ldots, a_{n})\) . Since \(f^{2} = f\) , one has \(j^{2} = j\) . By cor. 1 of I, p.~66, we have \(\chi(j) = 1\) if \(\chi \in \mathrm{U}_{1}\) and \(\chi(j) = 0\) if \(\chi \in \mathrm{U}_{2}\) , which proves the existence of the required idempotent.

On the other hand, if \(j_{1}\) is an idempotent of A, the relations \(j^{2} = j\) and \(j_{1}^{2} = j_{1}\) imply \((j - j_{1})(j + j_{1} - 1) = 0\) . If \(\mathcal{G}(j_{1}) = \mathcal{G}(j)\) , the Gelfand transform of \(j + j_{1} - 1\) takes values in \(\{-1, 1\}\) , hence \(j + j_{1} - 1\) is invertible (prop. 6 of I, p.~37), whence \(j = j_{1}\) .

COROLLARY. --- Let A be a commutative unital Banach algebra. The following assertions are equivalent:

\begin{enumerate}
\def\labelenumi{\alph{enumi})}
\item
  The space of characters X(A) is not connected;
\item
  There exists an idempotent element of A different from 0 and 1.
\item
  The algebra A is isomorphic to the product of two nonzero Banach algebras.
\end{enumerate}

The proposition demonstrates that \(a\) ) implies \(b\) ). If \(j\) is an idempotent of A, let \(\mathrm{I}_1 = j\mathrm{A}\) and \(\mathrm{I}_2 = (1 - j)\mathrm{A}\) . Then \(\mathrm{I}_1\) and \(\mathrm{I}_2\) are closed ideals of A, and \(\mathrm{I}_1 + \mathrm{I}_2 = \mathrm{A}\) . If \(j \notin \{0,1\}\) , the ideals \(\mathrm{I}_1\) and \(\mathrm{I}_2\) are distinct from A. On the other hand, the ideal \(\mathrm{I}_1\) (resp. \(\mathrm{I}_2\) ) is the set of elements \(x\) of A such that \(jx = x\) (resp. \((1 - j)x = x\) ), hence \(\mathrm{I}_1 \cap \mathrm{I}_2 = \{0\}\) . The algebra A is then identified with the product \(\mathrm{A} / \mathrm{I}_1 \times \mathrm{A} / \mathrm{I}_2\) . Thus, the assertion \(b\) ) implies \(c\) ). Finally, if A is isomorphic to \(\mathrm{A}_1 \times \mathrm{A}_2\) , the space \(\mathsf{X}(\mathsf{A})\) is identified with the sum space of \(\mathsf{X}(\mathsf{A}_1)\) and of \(\mathsf{X}(\mathsf{A}_2)\) (I, p.~6, no. 6), hence \(c\) ) implies \(a\) ).

PROPOSITION 14. --- Let A be a commutative Banach algebra without radical. For A to admit an identity element, it is necessary and sufficient that X(A) be compact.

The condition is necessary (I, p.~29, corollary). Suppose \(\mathsf{X}(\mathsf{A})\) compact. Let \(\widetilde{\mathsf{A}}\) be the Banach algebra obtained from A by adjoining an identity element, and identify \(\mathsf{X}'(\mathsf{A})\) with \(\mathsf{X}(\widetilde{\mathsf{A}})\) . The complement of \(\mathsf{X}(\mathsf{A})\) in \(\mathsf{X}(\widetilde{\mathsf{A}})\) is reduced to the character \(\chi_0\) of \(\widetilde{\mathsf{A}}\) whose kernel is A. The sets \(\mathsf{X}(\mathsf{A})\) and \(\{\chi_0\}\) are open in \(\mathsf{X}(\widetilde{\mathsf{A}})\) . By prop. 13, there exists an element \(j \in \mathsf{A}\) such that \(\chi(j) = 1\) for \(\chi \in \mathsf{X}(\mathsf{A})\) , and \(\chi_0(j) = 0\) . We therefore have \(j \in \mathsf{A}\) .

Let then \(x\) in A. We have \(\chi(jx) = \chi(x)\) for all \(\chi \in \mathsf{X}(\mathsf{A})\) , hence \(jx = x\) since A is without radical. Thus, \(j\) is an identity element of A.

PROPOSITION 15. --- Let A be a commutative Banach algebra, let \(I_{1}\) be an ideal of A and \(F_{1}\) be the set of \(\chi \in \mathsf{X}(A)\) which vanish on \(I_{1}\) . Let \(F_{2}\) be a subset of X(A) disjoint from \(F_{1}\) , closed for the Jacobson topology, and compact for the weak topology. Then there exists \(u \in \text{I}_{1}\) such that \(\mathcal{G}(u) = 1\) on \(F_{2}\) .

Let \(I_{2}\) be the intersection of the kernels of the characters belonging to \(F_{2}\) . The Banach algebra A/ \(I_{2}\) is without radical (prop. 8 of I, p.~38). Since \(F_{2}\) is closed for the Jacobson topology, the only elements of X(A) vanishing on \(I_{2}\) are those of \(F_{2}\) (cf.~I, p.~13). Hence \(F_{2}\) , equipped with the topology induced by the weak topology of X(A), is identified with X(A/ \(I_{2}\) ) equipped with the weak topology (I, p.~9, no. 7). Since \(F_{2}\) is weakly compact, the algebra A/ \(I_{2}\) possesses an identity element (prop. 14).

We then have \(I_{1} + I_{2} = A\) . Indeed, in the contrary case, \((I_{1} + I_{2})/I_{2}\) would be a proper ideal, hence contained in the kernel of a nonzero character of \(A/I_{2}\) (I, p.~30, th. 2). This character would define, by composition with the canonical projection \(A \rightarrow A/I_{2}\) , a nonzero character \(\chi\) of A vanishing on \(I_{1}\) and \(I_{2}\) , and would therefore belong to \(F_{1} \cap F_{2}\) , contrary to the hypothesis.

Since \(I_{1} + I_{2} = A\) , there exists \(u \in I_{1}\) whose class in \(A/I_{2}\) is an identity element of \(A/I_{2}\) . Then \(\chi(u) = 1\) for all \(\chi \in F_{2}\) , which concludes the proof.

COROLLARY. --- Let A be a commutative Banach algebra. Let \(F_{1}\) and \(F_{2}\) two disjoint subsets of \(\mathsf{X}(\mathsf{A})\) , closed for the Jacobson topology. We assume \(F_{2}\) weakly compact. Then there exists \(u \in A\) such that \(\mathcal{G}(u) = 1\) on \(F_{2}\) and \(\mathcal{G}(u) = 0\) on \(F_{1}\) .

\subsection{Partitions of the spectrum of an element}

Let A be a unital Banach algebra, \(x \in A\) , and \(K = \operatorname{Sp}_{\mathrm{A}}(x)\) . Let \(\Pi\) denote the set of subsets of K that are both open and closed in K. Let B be the closed full subalgebra of A generated by x; it is commutative.

For every \(H \in \Pi\) , there exists a unique element \(f_{H}\) of \(\mathcal{O}(K)\) equal to 1 in a neighborhood of H and to 0 in a neighborhood of K---H. We set \(j_{\mathrm{H}} = f_{\mathrm{H}}(x)\) . The element \(j_{H}\) is an idempotent of A, called associated with x and H, and we have the following formulas:

\[
j _ {\mathrm{H} \cap \mathrm{H} ^ {\prime}} = j _ {\mathrm{H}} j _ {\mathrm{H} ^ {\prime}} = j _ {\mathrm{H} ^ {\prime}} j _ {\mathrm{H}}\tag{15}
\]

(H, H' ∈ Π)

\[
j _ {\mathrm {H\cup H^ {\prime}}} = j _ {\mathrm{H}} + j _ {\mathrm {H^ {\prime}}} - j _ {\mathrm {H^ {\prime}}} j _ {\mathrm{H}}\tag{16}
\]

\[
(\mathrm{H}, \mathrm{H} ^ {\prime} \in \Pi)
\]

\[
j _ {\emptyset} = 0, \qquad j _ {\mathrm{K}} = 1.
\]

Let \(H \in \Pi\). We define \(A_{H} = j_{H}Aj_{H}\). This is a closed subalgebra of A, admitting the identity element \(j_{H}\) (cf.~Lemma 1 of I, p.~2). One also sets \(B_{H} = j_{H}Bj_{H}\) and \(x_{H} = xj_{H} = j_{H}x = j_{H}xj_{H} \in B_{H}\) .

Let \(g_{H}\) be the element of \(\mathcal{O}(K)\) defined by \(g_{\mathrm{H}}(z)=z\) in the neighborhood of H and \(g_{\mathrm{H}}(z)=0\) in the neighborhood of K---H. We have \(g_{\mathrm{H}}(z)=f_{\mathrm{H}}(z)z\) on K, and therefore \(x_{\mathrm{H}}=g_{\mathrm{H}}(x)\). It follows that, if \(H\neq K\) , one has

\[
\mathrm{Sp} _ {\mathrm{A}} (x _ {\mathrm{H}}) = g _ {\mathrm{H}} (\mathrm{K}) = \mathrm{H} \cup \{0 \}.
\]

Let \(\lambda \in \mathbf{C} - \mathrm{H}\) . Let \(h_{\mathrm{H},\lambda}\) be the element of \(\mathcal{O}(\mathrm{K})\) equal to \((\lambda - z)^{-1}\) in a neighborhood of H and to 0 in a neighborhood of K-H. We have \(h_{\mathrm{H},\lambda} = f_{\mathrm{H}}h_{\mathrm{H},\lambda}\) and \((\lambda f_{\mathrm{H}} - g_{\mathrm{H}})h_{\mathrm{H},\lambda} = f_{\mathrm{H}}\) . If we denote \(\mathrm{R_H}(x,\lambda) = h_{\mathrm{H},\lambda}(x)\) , we therefore have \(\mathrm{R_H}(x,\lambda) \in \mathrm{B_H}\) and

\[
\begin{array}{r} \mathrm {R_ {H}} (x, \lambda) (\lambda j _ {\mathrm{H}} - x _ {\mathrm{H}}) = (\lambda j _ {\mathrm{H}} - x _ {\mathrm{H}}) \mathrm {R_ {H}} (x, \lambda) = j _ {\mathrm{H}}, \\ \mathrm {R_ {H}} (x, \lambda) j _ {\mathrm {K-H}} = j _ {\mathrm {K-H}} \mathrm {R_ {H}} (x, \lambda) = 0. \end{array}\tag{17}
\]

In particular, \(\lambda \in \mathbf{C} - \mathrm{Sp}_{\mathrm{A_H}}(x_\mathrm{H})\) .

Now let \(\lambda\in\mathrm{H}\) . Suppose that \(\lambda j_{\mathrm{H}}-x_{\mathrm{H}}\) admits an inverse \(y\) in \(\mathrm{A}_{\mathrm{H}}\) . Using the formulas \(j_{\mathrm{H}}y=y\) (since \(y\in\mathrm{A}_{\mathrm{H}}\) ) and \(j_{\mathrm{K-H}}\mathrm{R}_{\mathrm{K-H}}(x,\lambda)=\mathrm{R}_{\mathrm{K-H}}(x,\lambda)\) (since \(\mathrm{R}_{\mathrm{K-H}}(x,\lambda)\in\mathrm{A}_{\mathrm{K-H}}\) ), we find

\[
(\lambda - x) (y + \mathrm{R} _ {\mathrm{K-H}} (x, \lambda)) = (\lambda - x) (j _ {\mathrm{H}} y + j _ {\mathrm{K-H}} \mathrm{R} _ {\mathrm{K-H}} (x, \lambda)) =
\]

\[
(\lambda j _ {\mathrm{H}} - x j _ {\mathrm{H}}) y + (\lambda j _ {\mathrm {K-H}} - x j _ {\mathrm {K-H}}) \mathrm{R} _ {\mathrm {K-H}} (x, \lambda) = j _ {\mathrm{H}} + j _ {\mathrm {K-H}} = 1,
\]

(thanks to formula (17) applied to K---H). One checks similarly that

\[
(y + \mathrm{R} _ {\mathrm{K-H}} (x, \lambda)) (\lambda - x) = 1.
\]

This proves that \(\lambda - x\) admits in A the inverse \(y + \mathrm{R}_{\mathrm{K - H}}(x, \lambda)\) , which is absurd. Thus we have \(\lambda \in \mathrm{Sp}_{\mathrm{A_H}}(x_{\mathrm{H}})\) . We therefore conclude that

\[
\mathrm{Sp} _ {\mathrm{A} _ {\mathrm{H}}} (x _ {\mathrm{H}}) = \mathrm{H}.\tag{18}
\]

In particular, if H is nonempty, the idempotent \(j_{H}\) is nonzero.

Formulas (17) and (18) prove that the function \(\lambda \mapsto \mathrm{R}_{\mathrm{H}}(x,\lambda)\) , defined in \(\mathbf{C} - \mathbf{H}\) , is the resolvent of \(x_{\mathrm{H}}\) with respect to \(\mathrm{A_H}\) .

PROPOSITION 16. --- We keep the previous notation. Let \((\mathrm{H}_{i})_{1\leqslant i\leqslant n}\) be a partition of \(\mathrm{Sp}_{\mathrm{A}}(x)\) into elements of \(\Pi\) .

\begin{enumerate}
\def\labelenumi{\alph{enumi})}
\item
  The algebra B is canonically identified with the algebra \(\mathrm{B_{H_1}}\times \dots \times \mathrm{B_{H_n}}\)
\item
  We have \(x_{\mathrm{H}_i}x_{\mathrm{H}_j} = 0\) for \(i\neq j\) , and
\end{enumerate}

\[
x = x _ {\mathrm{H} _ {1}} + x _ {\mathrm{H} _ {2}} + \dots + x _ {\mathrm{H} _ {n}};
\]

\begin{enumerate}
\def\labelenumi{\alph{enumi})}
\setcounter{enumi}{2}
\tightlist
\item
  We have
\end{enumerate}

\[
\mathrm{R} (x, \lambda) = \mathrm{R} _ {\mathrm{H} _ {1}} (x, \lambda) + \dots + \mathrm{R} _ {\mathrm{H} _ {n}} (x, \lambda)\tag{19}
\]

for all \(\lambda\in\mathbf{C}-\mathrm{Sp}_{\mathrm{A}}(x)\) . In particular, if \(\mathrm{H}\in\Pi\) , the resolvent \(\lambda\mapsto\mathrm{R}(x,\lambda)\) is equal in the neighborhood of \(\mathrm{H}\) to the sum of \(\mathrm{R}_{\mathrm{H}}(x,\lambda)\) and of a holomorphic function.

The relation \(1 = j_{H_{1}} + \cdots + j_{H_{n}}\) is a decomposition of 1 into pairwise orthogonal idempotents of B, so the algebra B is canonically identified with the product algebra \(B_{H_{1}} \times \cdots \times B_{H_{n}}\) (A, I, p.~105, prop. 10).

Assertion \(b\) ) follows from the corresponding relations for the functions \(g_{\mathrm{H}_i}\) ; assertion \(c\) ) is a consequence of \(a\) ) and of the equality \(\mathrm{R}(x_{\mathrm{H}},\lambda) = \mathrm{R}_{\mathrm{H}}(x,\lambda)\) .

PROPOSITION 17. --- Let \(\mu\) be an isolated point of \(\mathrm{Sp}_{\mathrm{A}}(x)\) . Then

\begin{enumerate}
\def\labelenumi{\alph{enumi})}
\tightlist
\item
  For every \(\lambda \in \mathbf{C} - \mathrm{Sp}_{\mathrm{A}}(x)\) , we have
\end{enumerate}

\[
\mathrm{R} (x, \lambda) = \mathrm{R} _ {\{\mu \}} (x, \lambda) + \mathrm{R} _ {\mathrm{Sp} _ {\mathrm{A}} (x) - \{\mu \}} (x, \lambda);
\]

\begin{enumerate}
\def\labelenumi{\alph{enumi})}
\setcounter{enumi}{1}
\item
  The function \(\lambda \mapsto \mathrm{R}_{\mathrm{Sp}_{\mathrm{A}}(x)-\{\mu\}}(x,\lambda)\) is holomorphic in \(\mathbf{C}-\mathrm{Sp}_{\mathrm{A}}(x)\) and in a neighborhood of \(\mu\); moreover, the function \(\lambda \mapsto \mathrm{R}_{\{\mu\}}(x,\lambda)\) is holomorphic in \(\mathbf{C}-\{\mu\}\) ;
\item
  We have
\end{enumerate}

\[
\lim _ {n \to + \infty} \| (x - \mu) ^ {n} j _ {\{\mu \}} \| ^ {1 / n} = 0
\]

and, for \(\lambda \in \mathbf{C} - \{\mu\}\) , the formula

\[
\mathrm{R} _ {\{\mu \}} (x, \lambda) = \sum_ {n = 0} ^ {\infty} (\lambda - \mu) ^ {- n - 1} (x - \mu) ^ {n} j _ {\{\mu \}}.\tag{20}
\]

The preceding implies the assertions \(a\) ) and \(b\) ). Let us prove \(c\) ). By replacing \(x\) by \(x-\mu\) , one reduces to the case where \(\mu=0\) . Set \(H = \{0\}\); this is an open and closed subset of \(\mathrm{Sp}_{\mathrm{A}}(x)\) . By formula (18), the spectrum of \(x_{H}\) in \(\mathrm{A}_{H}\) is \{0\}, hence \(x_{H}\) is quasi-nilpotent, that is, \(\|x^{n}j_{H}\|^{1/n}=\|(xj_{H})^{n}\|^{1/n}\) tends to 0 as \(n\) tends to + \(\infty\) . Moreover, for \(\lambda\neq0\) , in A one has

\[
\left(\lambda j _ {\mathrm{H}} - x _ {\mathrm{H}}\right) ^ {- 1} = \sum_ {n = 0} ^ {\infty} \lambda^ {- n - 1} x _ {\mathrm{H}} ^ {n}
\]

(Theorem 1 of I, p.~24, d)), whence (20).

COROLLARY 1. --- Let \(\mu\) be an isolated point of \(\mathrm{Sp}_{\mathrm{A}}(x)\) and \(p\) be a strictly positive integer. For \(\mu\) to be a pole of order \(p\) of the resolvent of \(x\) (cf.~VAR, R1, p.~30, 3.3.9), it is necessary and sufficient that \((x - \mu)^{p-1} j_{\{\mu\}} \neq 0\) and \((x - \mu)^{p} j_{\{\mu\}} = 0\) .

COROLLARY 2. --- Let \(\mu\) be an isolated point of \(\mathrm{Sp}_{\mathrm{A}}(x)\) . Let \(\Gamma\) be the oriented boundary of an open disk \(\Delta\) with center \(\mu\) such that

\[
\operatorname{Sp} _ {\mathrm{A}} (x) \cap (\Gamma \cup \Delta) = \{\mu \}.
\]

Then the idempotent \(j_{\{\mu \}}\) associated with \(x\) and \(\{\mu\}\) is given by

\[
j _ {\{\mu \}} = \frac {1}{2 i \pi} \int_ {\Gamma} (z - x) ^ {- 1} d z.
\]

*In other words, the idempotent \(j_{\{\mu\}}\) is the residue at \(\mu\) of the resolvent of \(x.*\)

For \(z\in \mathbf{C} - \mathrm{Sp}_{\mathrm{A}}(x)\) , one has

\[
(z - x) ^ {- 1} = \mathrm{R} (x, z) = \mathrm{R} _ {\{\mu \}} (x, z) + \mathrm{R} _ {\mathrm{H}} (x, z),
\]

where H = SpA(x) - \{μ\} (formula (19)). The function z ↦ R\_H(x, z) is holomorphic in C-H and in a neighborhood of \{μ\} (prop. 17, b)), hence

\[
\frac {1}{2 i \pi} \int_ {\Gamma} \mathrm{R} _ {\mathrm{H}} (x, z) d z = 0
\]

(VAR, R2, p.~48, 11.2.5). The function \(z \mapsto \mathrm{R}_{\{\mu\}}(x,z)\) is the resolvent of the element \(j_{\{\mu\}}xj_{\{\mu\}}\) of the unital algebra \(\mathrm{A}_{\{\mu\}}\) . We then have

\[
j _ {\{\mu \}} = \frac {1}{2 i \pi} \int_ {\Gamma} \mathrm{R} _ {\{\mu \}} (x, z) d z
\]

by prop. 9 of I, p.~75 applied to \(A_{\{\mu\}}\) and to the constant function 1 in a neighborhood of \(\Delta \cup \Gamma\) . The corollary follows.

\subsection{Holomorphic functional calculus in a complete normable real or complex algebra}

Let E be a real topological vector space. The topological vector space \(C \otimes E\), the complexification of E (EVT, II, p.~65), is denoted \(\mathrm{E}_{(\mathbf{C})}\) and E is identified with a real topological vector subspace of \(\mathrm{E}_{(\mathbf{C})}\) by the map \(x \mapsto 1 \otimes x\) .

PROPOSITION 18. --- The complex topological vector space \(E_{(\mathbf{C})}\) is normable (resp. complete) if and only if E is normable (resp. complete).

The underlying real topological vector space of \(\mathrm{E}_{(\mathbf{C})}\) is isomorphic to \(E \times E\) . Thus \(\mathrm{E}_{(\mathbf{C})}\) is complete if and only if E is complete, and E is normable if \(\mathrm{E}_{(\mathbf{C})}\) is.

Suppose conversely that E is normable. Let p be a norm defining the topology of E and B the unit ball of p. There exists a closed balanced neighborhood V of 0 in \(\mathrm{E}_{(\mathbf{C})}\) contained in \(B+iB\) (EVT, II, p.~66). The sets \(\lambda V\) , where \(\lambda\) describes \(R_{+}^{*}\) , therefore form a fundamental system of neighborhoods of 0 in \(\mathrm{E}_{(\mathbf{C})}\) . The gauge of V is a norm on \(\mathrm{E}_{(\mathbf{C})}\) that defines the topology of \(\mathrm{E}_{(\mathbf{C})}\) , hence \(\mathrm{E}_{(\mathbf{C})}\) is normable.

Remark. --- Let E and F be normable topological vector spaces over K. The vector space \(\mathcal{L}(\mathrm{E};\mathrm{F})\) of continuous linear maps from E into F, equipped with the topology of bounded convergence, is a normable topological vector space (EVT, III, p.~14).

Let E and F be normable topological vector spaces over R. The C-linear map \(\varphi\colon\mathcal{L}(\mathrm{E};\mathrm{F})_{(\mathbf{C})}\to\mathcal{L}(\mathrm{E}_{(\mathbf{C})};\mathrm{F}_{(\mathbf{C})})\) defined by \(\varphi(\lambda\otimes u)=\lambda u_{(\mathbf{C})}\) is an isomorphism of complex topological vector spaces. In particular, the dual of \(\mathrm{E}_{(\mathbf{C})}\) is identified with the complexification of the dual of E, and the normed algebra \(\mathcal{L}(\mathrm{E}_{(\mathbf{C})})\) with the complexification of the normed algebra \(\mathcal{L}(\mathrm{E})\) .

Let S be a compact subset of C stable under complex conjugation. Consider the C-algebra \(\mathcal{O}(S)\) of germs of complex-valued holomorphic functions in a neighborhood of S, endowed with the structure of a complex locally convex space defined in No. 1 of I, p.~49. If U is an open neighborhood of S in C, and \(h: U \to C\) be a holomorphic function, the image V of U under complex conjugation is an open neighborhood of S in C and \(h^{*}: w \mapsto \overline{h(\overline{w})}\) is a holomorphic function on V. Passing to the inductive limit yields a continuous involution \(f\mapsto f^{*}\) on the algebra \(\mathcal{O}(\mathrm{S})\) . In particular, we have:

\[
(f + g) ^ {*} = f ^ {*} + g ^ {*} \quad (f g) ^ {*} = f ^ {*} g ^ {*} \quad (\lambda f) ^ {*} = \bar {\lambda} f ^ {*}
\]

for \(f,g\) in \(\mathcal{O}(\mathrm{S})\) and \(\lambda\) in \(\mathbf{C}\) .

We denote by \(\mathcal{O}_{\mathbf{R}}(\mathrm{S})\) the set of germs \(f \in \mathcal{O}(\mathrm{S})\) such that \(f = f^{*}\) . It is a closed full sub- \(\mathbf{R}\) -algebra of \(\mathcal{O}(\mathrm{S})\) .

PROPOSITION 19. --- Let z denote the germ in \(\mathcal{O}(\mathrm{S})\) of the identity mapping of \(\mathbf{C}\) . Then \(\mathcal{O}_{\mathbf{R}}(\mathrm{S})\) is the smallest closed full sub- \(\mathbf{R}\) -algebra of \(\mathcal{O}(\mathrm{S})\) containing \(z\) .

We have \(z^{*} = z\) , hence z belongs to \(\mathcal{O}_{\mathbf{R}}(\mathrm{S})\). Let B be a closed full sub-R-algebra of \(\mathcal{O}(\mathrm{S})\) containing z. The map \(f \mapsto f + f^{*}\) of \(\mathcal{O}(\mathrm{S})\) into \(\mathcal{O}_{\mathbf{R}}(\mathrm{S})\) is continuous and surjective, and the set of germs of rational holomorphic functions in a neighborhood of S is dense in \(\mathcal{O}(\mathrm{S})\) (Theorem 3 of I, p.~69). To show that B contains \(\mathcal{O}_{\mathbf{R}}(\mathrm{S})\) it therefore suffices to show that if f is the germ of such a rational function, then we have \(f + f^{*} \in B\) .

There exist polynomials P and Q in \(\mathbf{C}[X]\) such that Q does not vanish at any point of S and such that we have \(f = \frac{\mathrm{P}(z)}{\mathrm{Q}(z)}\) . Denote by \(P^{*}\) and \(Q^{*}\) the polynomials obtained by replacing the coefficients of P and Q by their conjugates. We then have \(\mathrm{P}(z)^{*} = \mathrm{P}^{*}(z)\) and \(\mathrm{Q}(z)^{*} = \mathrm{Q}^{*}(z)\). Since S is stable under complex conjugation, the polynomial \(Q^{*}\) vanishes at no point of S. The germs \(\mathrm{Q}^{*}(z)\) and \((\mathrm{QQ}^{*})(z)\) are therefore invertible in \(\mathcal{O}(\mathrm{S})\) , and

\[
f + f ^ {*} = \frac {\mathrm{P} (z)}{\mathrm{Q} (z)} + \frac {\mathrm{P} ^ {*} (z)}{\mathrm{Q} ^ {*} (z)} = \frac {(\mathrm{PQ} ^ {*} + \mathrm{P} ^ {*} \mathrm{Q}) (z)}{(\mathrm{QQ} ^ {*}) (z)}.
\]

Since the polynomials PQ* + P*Q and QQ* have real coefficients and B is a full R-subalgebra of \(\mathcal{O}(\mathrm{S})\) containing \(z\) , the element \(f + f^{*}\) belongs to B. This concludes the proof of the proposition.

Let A be a complete normable unital algebra over R. Let x be an element of A. The spectrum of the element \(1 \otimes x\) of the algebra \(\mathrm{A}_{(\mathbf{C})}\) is called the complex spectrum of x, and is denoted \(\mathrm{Sp}_{\mathrm{A}_{(\mathbf{C})}}(x)\) . Its intersection with the set R is none other than the spectrum \(\mathrm{Sp}_{\mathrm{A}}(x)\) of x relative to A, which is sometimes called the real spectrum of x. The complex spectrum \(\mathrm{Sp}_{\mathrm{A}_{(\mathbf{C})}}(x)\) is a compact subset of C, stable under complex conjugation. It is nonempty when the algebra A is not reduced to 0.

Let \(x\) be an element of A. The spectral radius of \(1 \otimes x \in \mathrm{A}_{(\mathbf{C})}\) is equal to the spectral radius \(\varrho(x)\) of \(x\) . It is the smallest real number \(r \geqslant 0\) such that \(|\lambda| \leqslant r\) for all \(\lambda \in \mathrm{Sp}_{\mathrm{A}_{(\mathbf{C})}}(x)\) . We have

\[
\varrho (x) = \lim _ {n \to + \infty} \| x ^ {n} \| ^ {1 / n} = \inf _ {n > 0} \| x ^ {n} \| ^ {1 / n}
\]

for every norm on A that defines the topology of A. Indeed, one may assume that the norm on A is the restriction of a norm on \(\mathrm{A}_{(\mathbf{C})}\) that defines the topology of \(\mathrm{A}_{(\mathbf{C})}\) and apply prop. 1 of I, p.~20.

We denote by \(u \mapsto \overline{u}\) the endomorphism of the R-algebra \(A_{(C)}\) which maps \(\lambda \otimes a\) to \(\overline{\lambda} \otimes a\) . It is continuous.

Lemma. --- For every \(f \in \mathcal{O}(\mathrm{Sp}_{\mathrm{A}_{(\mathbf{C})}}(x))\) , one has \(f^{*}(1 \otimes x) = \overline{f(1 \otimes x)}\) .

The maps \(f \mapsto f(1 \otimes x)\) and \(f \mapsto \overline{f^*(1 \otimes x)}\) are continuous unital homomorphisms of \(\mathbf{C}\) -algebras from \(\mathcal{O}(\mathrm{Sp}(x))\) into \(\mathrm{A}_{(\mathbf{C})}\) which map \(z\) to \(1 \otimes x\); therefore they are equal (I, p.~74, th. 5).

PROPOSITION 20. --- For every \(f \in \mathcal{O}_{\mathbf{R}}(\mathrm{Sp}_{\mathrm{A}_{(\mathbf{C})}}(x))\) , there exists a unique element \(f(x)\) of A such that \(f(1 \otimes x) = 1 \otimes f(x)\) in \(\mathrm{A}_{(\mathbf{C})}\) . The map \(f \mapsto f(x)\) of \(\mathcal{O}_{\mathbf{R}}(\mathrm{Sp}_{\mathrm{A}_{(\mathbf{C})}}(x))\) in A is the unique continuous unital homomorphism of R-algebras that maps the germ in \(\mathcal{O}_{\mathbf{R}}(\mathrm{Sp}_{\mathrm{A}_{(\mathbf{C})}}(x))\) of the identity map of C to \(x\).

We denote \(S = \text{Sp}_{\text{A}(\mathbf{C})}(x)\) . By the above lemma, for every germ \(f \in \mathcal{O}_{\mathbf{R}}(\text{Sp}(x))\) , one has \(f(1 \otimes x) = \overline{f(1 \otimes x)}\). The first assertion follows from this. Let us denote by \(z\) the germ in \(\mathcal{O}_{\mathbf{R}}(\text{S})\) of the identity mapping of \(\mathbf{C}\) . The map \(f \mapsto f(x)\) is a continuous unital homomorphism of the \(\mathbf{R}\) -algebra \(\mathcal{O}_{\mathbf{R}}(\text{Sp}(x))\) into \(A\) , which maps \(z\) to \(x\). This is the unique one by Prop. 19, since any morphism with these properties is uniquely determined on every sub- \(\mathbf{R}\) -algebra of \(\mathcal{O}(\text{S})\) containing \(z\) .

Let \(f \in \mathcal{O}_{\mathbf{R}}(\mathrm{Sp}_{\mathrm{A}(\mathbf{C})}(x))\) . The element \(f(x)\) belongs to every closed full subalgebra of A containing \(x\) (prop. 19), hence belongs to the bicommutant of \(x\) in A. The complex spectrum of \(f(x)\) is equal to \(f(\mathrm{Sp}(x))\) (I, p.~75, prop. 8). For all \(g \in \mathcal{O}_{\mathbf{R}}(f(\mathrm{Sp}_{\mathrm{A}(\mathbf{C})}(x)))\) , one has \(g \circ f \in \mathcal{O}_{\mathbf{R}}(\mathrm{Sp}_{\mathrm{A}(\mathbf{C})}(x))\) and (loc. cit.) \(g \circ f(x) = g(f(x))\) .

Let U be an open subset of C, stable under complex conjugation. The set \(\Omega\) of elements x of A whose complex spectrum is contained in U is open in A (I, p.~76, prop. 10). Let f be a holomorphic function on U such that \(f^{*} = f\) . The map \(x \mapsto f(x)\) of \(\Omega\) in A is analytic (loc. cit.).

Let A, B be complete normable unital associative algebras over R and \(\varphi\colon A\to B\) a continuous unital algebra morphism. Let \(x\in A\) . The complex spectrum of \(\varphi(x)\) is contained in that of x and, for every \(f\in\mathcal{O}_{\mathbf{R}}(\mathrm{Sp}_{\mathrm{A}(\mathbf{C})}(x))\) , one has \(f(\varphi(x))=\varphi(f(x))\) . This follows immediately from the analogous statement in the complex case (I, p.~75, prop. 8).

\subsection{Case of an algebra without unit element}

Let A be a complete normable algebra, not necessarily unital, over K = R or C. Denote by ( \(\widetilde{A}, e\) ) the unital algebra obtained from A by adjoining an identity element. It is normable and complete.

Let \(x\) be an element of A. If K = C, denote by \(\mathrm{Sp}'(x) = \mathrm{Sp}_{\widetilde{\mathrm{A}}} (x)\) the spectrum of \(x\) with respect to \(\widetilde{\mathrm{A}}\) , and consider a germ \(f \in \mathcal{O}(\mathrm{Sp}'(x))\) . If K = R, denote by \(\mathrm{Sp}'(x)\) the complex spectrum of the element \(x\) of \(\widetilde{\mathrm{A}}\) , and consider a germ \(f \in \mathcal{O}_{\mathbf{R}}(\mathrm{Sp}'(x))\). In these two cases, 0 belongs to \(\mathrm{Sp}'(x)\) , and the element \(f(x)\) of \(\widetilde{\mathrm{A}}\) belongs to A if and only if \(f\) satisfies \(f(0) = 0\) . Indeed, the projection \(\pi: \widetilde{\mathrm{A}} \to \mathrm{Ke}\) is a continuous morphism whose kernel is A, and we have \(\pi(f(x)) = f(\pi(x)) = f(0)\) .

\section{REGULAR COMMUTATIVE BANACH ALGEBRAS}

In this section, the base field is C.

\subsection{Definition}

PROPOSITION 1. --- Let A be a commutative Banach algebra. The following conditions are equivalent:

\begin{enumerate}
\def\labelenumi{(\roman{enumi})}
\item
  The weak topology and the Jacobson topology on \(\mathsf{X}(\mathsf{A})\) coincide;
\item
  For every \(\chi \in \mathsf{X}(\mathrm{A})\) and every weakly closed subset \(\mathrm{F}\) of \(\mathsf{X}(\mathrm{A})\) such that \(\chi \notin \mathrm{F}\) , there exists an \(x \in \mathrm{A}\) such that \(\mathcal{G}(x)\) is equal to 1 at \(\chi\) and to 0 on \(\mathrm{F}\) ;
\item
  For every weakly compact subset K and every weakly closed subset F of X(A) such that K ∩ F = ∅, there exists an element \(x \in A\) such that \(\mathcal{G}(x)\) is equal to 1 on K and to 0 on F.
\end{enumerate}

Let \(M \subset \mathsf{X}(A)\) . Saying that M is closed for the Jacobson topology means that, for every \(\chi \in \mathsf{X}(A) - \mathsf{M}\) , there exists an \(x \in A\) such that \(\mathcal{G}(x)\) vanishes on M but not at \(\chi\) (lemma 2 of I, p.~39). Condition (ii) therefore means that every subset of \(\mathsf{X}(A)\) weakly closed is closed for the Jacobson topology, which shows that (ii) \(\Longrightarrow\) (i). Moreover, (iii) \(\Longrightarrow\) (ii) since the set \(\{\chi\}\) is weakly compact in \(\mathsf{X}(A)\). Finally, (i) \(\Longrightarrow\) (iii) by the corollary of prop. 15 of I, p.~81.

DEFINITION 1. --- Let A be a commutative Banach algebra. It is said to be regular if it satisfies the equivalent conditions of Proposition 1.

Note. --- Let \(\tilde{\mathsf{A}}\) be the Banach algebra obtained from A by adjoining an identity element \(e\) . Condition (ii) of Prop. 1 shows that if \(\tilde{\mathsf{A}}\) is regular, then A is regular. Suppose A is regular and let us show that \(\tilde{\mathsf{A}}\) is regular. Consider subsets F and F′ of \(\mathsf{X}(\tilde{\mathsf{A}})\) which are disjoint and weakly closed (hence weakly compact), and let us construct an \(x \in \tilde{\mathsf{A}}\) such that \(\mathcal{G}(x)\) vanishes on F, and is equal to 1 on F′. Let \(\chi_0 \in \mathsf{X}(\tilde{\mathsf{A}})\) be the zero character on A. If \(\chi_0 \notin \mathrm{F}'\) , by condition (iii) of prop. 1 and the hypothesis on A, there exists an element \(x \in \mathsf{A}\) such that \(\mathcal{G}(x)\) vanishes on F and equals 1 on F'. If \(\chi_0 \in \mathrm{F}'\) , one has \(\chi_0 \notin \mathrm{F}\) ; there therefore exists an element \(y \in \mathsf{A}\) such that \(\mathcal{G}(y)\) vanishes on F' and equals 1 on F. The element \(x = e - y\) of \(\tilde{\mathsf{A}}\) then has the required property.

Examples. --- Let us return to the examples of no. 2 of I, p.~17.

The algebra of continuous complex-valued functions tending to 0 at infinity on a locally compact space X (example 3 of I, p.~17) is regular (cf.~I, p.~36, example 1).

The algebra of functions \(n\) times differentiable on \([0, 1]\) (example 4 of I, p.~18) is regular (cf.~I, p.~36, example 2).

If G is a commutative locally compact group and \(\mu\) be a Haar measure on G, then the algebra \(L^{1}(G,\mu)\) (example 7 of I, p.~19) is regular (cf.~II, p.~219, cor. 2).

The algebra of functions that are continuous in the disk \(|z| \leqslant 1\) and analytic in the interior (example 9 of I, p.~20) is not regular (cf.~I, p.~193, exerc. 6).

PROPOSITION 2. --- Let A be a commutative unital regular Banach algebra. Let \(n \geqslant 1\) be an integer and \((\mathrm{U}_{1}, \ldots, \mathrm{U}_{n})\) be an open covering of \(\mathsf{X}(\mathsf{A})\) . There exist elements \(x_{1}, \ldots, x_{n}\) of A with sum 1 such that \(\operatorname{Supp}(\mathcal{G}(x_{i})) \subset \mathrm{U}_{i}\) for \(i = 1, \ldots, n\) .

We prove the proposition by induction on n. The assertion is valid for n = 1. Suppose that \(n \geqslant 2\) and that the assertion is established for n - 1.

There exists an open covering \((\mathrm{V}_{1},\ldots,\mathrm{V}_{n})\) of \(\mathsf{X}(\mathsf{A})\) such that \(\overline{V}_{i}\subset U_{i}\) for every i. By the induction hypothesis, there exist elements \(x,x_{3},\ldots,x_{n}\in A\) such that \(x+x_{3}+\cdots+x_{n}=1\) and \(\operatorname{Supp}(\mathcal{G}(x))\subset\mathrm{V}_{1}\cup\mathrm{V}_{2}\) , \(\operatorname{Supp}(\mathcal{G}(x_{i}))\subset\mathrm{V}_{i}\) for \(i\geqslant3\) . Let \(\mathrm{K}=\operatorname{Supp}(\mathcal{G}(x))\subset\mathrm{V}_{1}\cup\mathrm{V}_{2}\) . Let \(K_{1}\) (resp. \(K_{2}\) ) be the set of elements of K that do not belong to \(V_{1}\) (resp. \(V_{2}\) ). Then \(K_{1}\) and \(K_{2}\) are disjoint compact subsets of K. Since the Banach algebra A is regular, there therefore exists \(y\in A\) such that \(\mathcal{G}(y)=1\) on \(K_{1}\) and \(\mathcal{G}(y)=0\) on \(K_{2}\) . Then \(\mathcal{G}(xy)\) is zero on \(\mathsf{X}(\mathsf{A})-\mathsf{K}\) and on \(K_{2}\) , hence \(\operatorname{Supp}\mathcal{G}(xy)\subset\overline{\mathrm{V}}_{2}\subset\mathrm{U}_{2}\) . Likewise, \(\mathcal{G}(x(1-y))\) is zero on \(\mathsf{X}(\mathsf{A})-\mathsf{K}\) and on \(K_{1}\) , hence \(\operatorname{Supp}\mathcal{G}(x(1-y))\subset\overline{\mathrm{V}}_{1}\subset\mathrm{U}_{1}\) . The elements \(x_{1}=x(1-y)\) , \(x_{2}=xy\) , and \(x_{3},\ldots,x_{n}\) then satisfy the properties of the proposition.

COROLLARY 1. --- Let A be a commutative unital regular Banach algebra, let I be an ideal of A, and let \(f: \mathsf{X}(\mathsf{A}) \to \mathbf{C}\) be a continuous function. Suppose that, for every \(\chi \in \mathsf{X}(\mathsf{A})\) , there exists an element \(y_{\chi} \in \mathrm{I}\) such that \(f = \mathcal{G}(y_{\chi})\) in a neighborhood of \(\chi\) . Then there exists an element \(y \in \mathrm{I}\) such that \(f = \mathcal{G}(y)\) .

Since X(A) is compact, there exists a finite open covering \((\mathrm{U}_{1},\ldots,\mathrm{U}_{n})\) of X(A), and elements \(y_{1},\ldots,y_{n}\) of I such that \(f=\mathcal{G}(y_{i})\) on \(U_{i}\) . By prop. 2, there exist elements \(x_{1},\ldots,x_{n}\) of A with sum 1 such that \(\operatorname{Supp}(\mathcal{G}(x_{i}))\subset\mathrm{U}_{i}\) for every i. Let \(y=x_{1}y_{1}+\cdots+x_{n}y_{n}\) . This is an element of I that has the required property. Indeed, let \(\chi\in\mathrm{X}(A)\) . For \(1\leqslant i\leqslant n\) , we have \(\mathcal{G}(x_{i})(\chi)\mathcal{G}(y_{i})(\chi)=\mathcal{G}(x_{i})(\chi)f(\chi)\) since \(\mathcal{G}(y_{i})(\chi)=f(\chi)\) if \(\chi\in U_{i}\) , and \(\mathcal{G}(x_{i})(\chi)=0\) if \(\chi\notin U_{i}\) . It follows that

\[
\mathcal {G} (y) (\chi) = \sum_ {i = 1} ^ {n} \mathcal {G} (x _ {i}) (\chi) \mathcal {G} (y _ {i}) (\chi) = f (\chi) \sum_ {i = 1} ^ {n} \mathcal {G} (x _ {i}) (\chi) = f (\chi).
\]

COROLLARY 2. --- Let A be a commutative regular Banach algebra, I an ideal of A, and \(f\colon\mathsf{X}'(\mathrm{A})\to\mathbf{C}\) be a continuous function. Suppose that, for every \(\chi\in\mathsf{X}'(\mathrm{A})\) , there exists an element \(y_{\chi}\in\mathrm{I}\) such that \(f=\mathcal{G}'(y_{\chi})\) in a neighborhood of \(\chi\) . Then there exists an element \(y\in\mathrm{I}\) such that \(f=\mathcal{G}'(y)\) .

Let \(\tilde{\mathsf{A}}\) be the Banach algebra obtained from \(\mathsf{A}\) by adjoining an identity element. Then \(\tilde{\mathsf{A}}\) is regular (remark 1), and \(\mathsf{X}'(\mathsf{A}) = \mathsf{X}(\widetilde{\mathsf{A}})\) ; it therefore suffices to apply cor. 1 to \(\tilde{\mathsf{A}}\) and to the ideal I.

If I is an ideal of a commutative Banach algebra, recall (cf.~I, p.~30) that we denote by V(I) the set of \(\chi \in \mathsf{X}(\mathrm{A})\) whose kernel contains I, in other words the set of \(\chi \in \mathsf{X}(\mathrm{A})\) where all the functions \(\mathcal{G}(x)\) for \(x \in \mathrm{I}\) vanish. This is a subset of \(\mathsf{X}(\mathrm{A})\) closed for the Jacobson topology.

PROPOSITION 3. --- Let A be a commutative regular Banach algebra, I an ideal of A, and K a compact subset of X(A) disjoint from V(I). There exists an element \(x \in I\) such that \(\mathcal{G}(x) = 1\) for every \(x\) in K.

This is a special case of prop. 15 of I, p.~81, taking into account the fact that the Jacobson topology coincides with the weak topology on X(A).
